    % =====================================================================
%  PAPER D  --  merged submission for NDSS (Network and Distributed
%  System Security Symposium), combining the cross-layer benchmark
%  (formerly Paper C) and the composition theorem (formerly Paper A).
%
%  NDSS FORMAT COMPLIANCE
%  ----------------------
%  NDSS requires: US Letter (not A4), two-column, columns no more than
%  9.25 in high and 3.5 in wide, Times font at 10 pt or larger with
%  11 pt or larger leading, and at most 13 pages excluding the Ethics
%  Considerations section, references, and appendices.
%
%  The class line below (IEEEtran, conference, letterpaper, 10pt, with
%  Times via mathptmx) reproduces exactly that geometry: \columnwidth
%  = 3.5 in and \textheight = 9.25 in (verified at compile time by the
%  \typeout block after \begin{document}).
%
%  FOR FINAL SUBMISSION: download the official template from
%  ndss-symposium.org/ndss2026/submissions/templates/ and replace the
%  two lines marked "NDSS SWAP" below with \documentclass{ndss}.
%  No change to the body is required.
% =====================================================================
\documentclass[conference,letterpaper,10pt]{IEEEtran}   % NDSS SWAP (1/2)
\IEEEoverridecommandlockouts
\usepackage{mathptmx}                                    % NDSS SWAP (2/2): Times
\usepackage{cite}
\usepackage{amsmath,amssymb,amsfonts}
\usepackage{amsthm}
\usepackage{graphicx}
\usepackage{booktabs}
\usepackage{tabularx}
\usepackage{subcaption}
\usepackage{xcolor}
\usepackage{tikz}
\usepackage{pgfplots}
\pgfplotsset{compat=1.17}
\usetikzlibrary{arrows.meta,positioning,fit,backgrounds,shapes.geometric,calc}
\usepackage{enumitem}
\usepackage[hidelinks]{hyperref}
\usepackage[htt]{hyphenat}

% NDSS caps column height at 9.25in; IEEEtran's default is 9.30in, so we
% trim it to be strictly within spec (9.25in = 668.5pt at 72.27pt/in).
\setlength{\textheight}{668.4pt}

% ---- anonymisation switch: NDSS review is double-blind ----------------
% \anontrue  = anonymised submission version (what NDSS requires)
% \anonfalse = named version for internal / supervisor review
\newif\ifanon
\anontrue
% Own work is cited anonymously in the review version; the named entry
% (uavewbench2026) returns with \anonfalse.
\newcommand{\ownbenchcite}{\ifanon\cite{uavewbench2026anon}\else\cite{uavewbench2026}\fi}

% ---- shared style tokens --------------------------------------------
\definecolor{netcol}{RGB}{31,78,121}
\definecolor{autocol}{RGB}{18,120,90}
\definecolor{bridgecol}{RGB}{176,96,0}
\definecolor{accent}{RGB}{140,20,20}
% Layer names are set as ordinary text. Colour is reserved for the figures,
% where it distinguishes the two layers visually; in running prose it is noise.
\newcommand{\layernet}{network}
\newcommand{\layerauto}{autonomy}
\newcommand{\dataset}[1]{\textsc{#1}}
\newcommand{\code}[1]{\texttt{#1}}
\newcommand{\Real}{\mathbb{R}}
\newcommand{\Norm}[1]{\left\lVert #1\right\rVert}
\newcommand{\MCR}{\mathrm{MCR}}

\newtheorem{theorem}{Theorem}
\newtheorem{lemma}{Lemma}
\newtheorem{definition}{Definition}
\newtheorem{proposition}{Proposition}

\newcommand{\fc}[1]{{\color{red}{FC: #1}}}

\begin{document}


% The title deliberately says CONDITIONAL and names the calibrated interface.
% An earlier version read "An End-to-End Framework for Network-to-Navigation UAV
% Security", which implies the entire chain is certified. It is not: the middle
% link is an empirical calibration, and the guarantee holds only under stated
% conditions. See \S\ref{sec:scope}.
\title{From Wire to Flight: A Conditional Mission-Completion Guarantee for UAVs}


\ifanon
  \author{\IEEEauthorblockN{Paper \#XXX}
  \IEEEauthorblockA{Submitted for double-blind review}}
\else
  \author{
  \IEEEauthorblockN{Roger Nick Anaedevha\IEEEauthorrefmark{1}, \textit{IEEE Member},
  Federico Cor\`o\IEEEauthorrefmark{2}, and Mauro Conti\IEEEauthorrefmark{2},
  \textit{IEEE Fellow}}
  \IEEEauthorblockA{\IEEEauthorrefmark{1}Institute of Cyber Intelligent Systems,
  National Research Nuclear University MEPhI, Moscow, Russia}
  \IEEEauthorblockA{\IEEEauthorrefmark{2}SPRITZ Security and Privacy Research Group,
  Department of Mathematics, University of Padua, Padua, Italy}}
\fi

\maketitle


\typeout{NDSS GEOMETRY CHECK: columnwidth=\the\columnwidth\space
(target 3.5in=252.94pt), textheight=\the\textheight\space
(target 9.25in=668.62pt)}

\begin{abstract}
Research on unmanned aerial vehicle (UAV) security proceeds along two lines that
do not intersect. Intrusion detection flags malicious traffic but says nothing
about the flight that follows a missed intrusion. Certified robustness bounds a
navigation model's response to a perturbation but does not ask where that
perturbation originates or how often an adversary can produce it. Neither
addresses the question an operator faces: whether a mission can still be
completed when the attacker remains just below the alarm threshold.

We address this question with a conditional guarantee. A two-layer schema over
six public datasets takes as its unit the \emph{cross-layer pairing}, a
network-attack interval joined to the navigation degradation it causes and
tagged with the strength of that causal evidence. We prove that, against an
attacker who delays relayed messages while remaining just below the threshold,
the proportion of completed missions has a lower bound computable from that
threshold. Two measurements make the bound informative. Hidden delay saturates
at $4.84$\,s across $15{,}392$ relay hops, and each second of delayed navigation
input causes at most $1.63$\,m of position error across $674$ post-onset samples
from real GPS-jamming and GPS-spoofing flights, an order of magnitude below the
kinematic worst case.

The chain is not proved end to end. Its middle link is calibrated on three hover
flights of a single airframe; the calibrated rate varies by $6.1\times$ with
attack type, satellite visibility and staleness, so the certificate uses its
supremum. The guarantee further assumes Lipschitz dynamics, estimator
re-anchoring within a stated interval, and calibration transfer, and it bounds a
population of missions rather than a single flight. Under these conditions the
composed defence certifies a non-zero floor, whereas neither single-layer
defence certifies one. The schema, adapters and certificate engine are released
as an open artifact.
\end{abstract}

% =====================================================================
\section{Introduction}
\label{sec:intro}
% =====================================================================
An autonomous UAV presents four distinct interfaces to an adversary. Commands
reach it from a ground station over a command-and-control (C2) radio link. Peer
aircraft in a swarm exchange position and status reports over a multi-hop mesh,
in which each aircraft relays traffic for its neighbours. Inside the airframe,
the flight controller drives motor controllers and reads sensors over a shared
bus. Finally, a navigation pipeline fuses satellite, inertial, and visual
measurements into the state estimate on which the flight control loop acts. A
large body of work defends each of these interfaces in isolation. We are
concerned with a question that none of these defences is designed to answer,
namely what happens to the flight once some attacks have got through.

The question has practical weight because the attacks of greatest operational
concern are not those that exfiltrate data. A jammer raises the noise floor at
the satellite-navigation receiver until the position solution
collapses~\cite{kassas2022jammer}. A compromised mesh relay forwards every
message it receives, but with a delay, so the aircraft acts on a peer's position
report that is already stale~\cite{datamut2026,lou2019timedelay}. In both cases
a network event degrades navigation, and the two lie at opposite ends of a
single causal chain. Neither existing datasets nor existing guarantees represent
this chain, so it cannot be analysed.

The practical consequence is a claim that neither research community can
currently substantiate. On the network side, a UAV intrusion detector can be
shown to flag delayed forwarding events or spoofed navigation messages at high
recall. What cannot be evaluated is the operationally decisive claim behind that
figure, namely that a detector tuned to a given false-alarm budget allows
aircraft to keep completing their missions; no public dataset records both a
network event and the flight it degrades. The navigation side has the converse
gap. Certified robustness establishes that a model's output is stable under
every input perturbation of magnitude at most $R$. This is a statement about a
model rather than a flight, and it says nothing about where a perturbation of
that magnitude originates or how often an adversary can produce one.

Closing this gap requires work at both ends simultaneously, since the two halves
depend on each other. A guarantee spanning both layers needs data objects that
record the coupling, and a corpus that records the coupling needs a guarantee to
show what those objects are for. We therefore unify six datasets, four covering
the network and two covering navigation, under a schema whose unit of analysis
is a \emph{cross-layer pairing} rather than a packet or a flight. A pairing is
an interval during which a network attack is active, joined to the interval of
degraded navigation it induces. On these objects we then prove that a detector's
alarm threshold determines a lower bound on the fraction of missions that still
complete, under the pessimistic assumption that the adversary knows the
threshold and stays below it. To our knowledge, this is the first guarantee to
carry a network-security threshold through to a mission-completion floor.

This claim requires an immediate qualification. Only one of the six datasets
observes a network event and the flight it degrades on the same airframe, and
its release does not include machine-readable attack intervals; consequently,
none of the pairings we release belongs to the strongest evidence class
(\S\ref{sec:provenance_dist}). The cross-layer coupling therefore rests on a
stated physical relation calibrated on real flights, rather than on observation
of both windows on the same platform. Accordingly, we claim a conditional
guarantee with an empirically calibrated interface (\S\ref{sec:scope}), not
mission-level security established end to end by measurement. The contributions
of the paper are as follows.

\subsection*{Contributions}
\begin{itemize}
\item \textbf{A two-layer schema, and an adapter for each of six datasets
(\S\ref{sec:schema},~\S\ref{sec:corpus}).} The schema provides a single
representation for any observation from any source, whether a network flow, a
bus frame, a telemetry sample, or an entire flight. An \emph{adapter} is the
program that reads one dataset in its native release format and emits records in
that representation; we write one per dataset and verify that no source column
is discarded. Each record also carries a flag indicating the strength of the
evidence for its cross-layer coupling, so that a reader can restrict any
analysis to the evidence classes they consider acceptable.

\item \textbf{A detector-operating-curve harness (\S\ref{sec:harness}).} A
detector's \emph{operating point} is the threshold at which it raises an alarm;
a low threshold catches more attacks but also produces more false alarms. We
sweep this threshold and measure recall and false-positive rate as functions of
it, which supplies the network-side premise required by a composed guarantee.
Across three topologies the curve exhibits a universal knee-and-floor shape, the
floor being the attacker's self-exposure probability.

\item \textbf{An empirically tightened interface, and a measurement of its
stability (\S\ref{sec:interface}).} We prove a residual-budget lemma and then
tighten it by measurement, showing that the undetected budget saturates at the
contact-window slack, which we verify at $\le 4.84$\,s over $15{,}392$ hops. We
calibrate the delay-to-position rate $\gamma$ on three real spoofing and jamming
flights, obtaining a bound an order of magnitude tighter than the kinematic
worst case. We then test whether $\gamma$ is constant and find that it is not.
Estimated at all $674$ post-onset samples rather than once per flight, it varies
by $6.1\times$, and it does so systematically, with attack type, satellite
visibility and staleness, each in a direction with a mechanistic explanation
(\S\ref{sec:gammastability}). The certificate is therefore evaluated at the
supremum over observed conditions rather than at an average, and the reported
spread is a lower bound on the true variability, since the released flights fix
platform and flight mode.

\item \textbf{A composition theorem and a four-certificate architecture
(\S\ref{sec:theorem},~\S\ref{sec:certs}).} We prove the composed
network-to-navigation guarantee and its soundness, and explain why the
deterministic core does not subsume randomised smoothing, PAC-Bayes, or
multiplicative-weights regret.

\item \textbf{Instantiation, ablation, and an open artifact
(\S\ref{sec:eval}).} Among the configurations we compare, the composed guarantee
alone yields a non-zero certified floor against an evading attacker, and under
the measured interface a deployed detector's operating point lies within the
certified window. The comparison is with ablations of our own construction
rather than with alternative assurance methods, for reasons given in
\S\ref{sec:dominance}. Every figure and table is generated from committed result
files, and every number carries a provenance tag.
\end{itemize}

\subsection*{Scope of the Guarantee}
\label{sec:scope}
Because a composed result is easily over-read, we state its boundaries here
rather than in the limitations section. We establish a conditional guarantee
under an empirically calibrated interface, and the guarantee is conditional in
three distinct senses.

\begin{enumerate}[label=(C\arabic*),leftmargin=2.6em]
\item\label{c:lip} \textbf{Dynamics.} The navigation drift is $L$-Lipschitz over
the operating region. Outside that region, for instance in a stall-recovery
transient or under actuator saturation, no finite $L$ bounds the drift, and the
certificate makes no statement.
\item\label{c:anchor} \textbf{Re-anchoring.} The estimator receives at least one
uncorrupted state update per certification horizon $T_c$. A delay adversary
satisfies this by construction; a persistent-denial adversary does not.
\item\label{c:interface} \textbf{Interface transfer.} The measured
delay-to-position rate $\gamma$ applies to the deployment in question. This is
the weakest of the three conditions: $\gamma$ is calibrated on three hover
flights of one airframe and, as \S\ref{sec:gammastability} shows, is not
constant. It varies by $6.1\times$ over the conditions we can observe, and we
cannot vary platform or flight mode at all.
\end{enumerate}

Two further boundaries should be stated. First, the chain
$\theta\mapsto\Delta\mapsto\delta\mapsto\MCR$ is not proved end to end. The
first and last links are proved; the middle link, $\Delta\mapsto\delta$, is an
empirical calibration. Describing the result as ``certified end-to-end
security'' would misrepresent where the evidence lies, and we avoid that
description. Second, MCR is a probability over a distribution of missions
(\S\ref{sec:probspace}); a bound on that proportion is not a bound for an
arbitrary individual mission drawn from a different distribution.

\emph{Organisation.}
\S\ref{sec:problem} states the threat model, defines mission completion, and
formalises the cross-layer pairing. \S\ref{sec:related} positions the work.
\S\ref{sec:schema} presents the two-layer schema and \S\ref{sec:corpus} the six
ingestion adapters; \S\ref{sec:harness} builds the detector-operating-curve
harness. \S\ref{sec:interface} derives and then measures the interface mapping,
\S\ref{sec:theorem} proves the composition theorem, and \S\ref{sec:certs}
places it within a four-certificate architecture. \S\ref{sec:eval} instantiates
and evaluates the result, and
\S\ref{sec:discussion}--\S\ref{sec:conclusion} present the discussion,
limitations, and conclusion.

% =====================================================================
\section{Threat Model and Problem Formulation}
\label{sec:problem}
% =====================================================================
This section fixes three elements on which the remainder of the paper depends:
the parts of the aircraft an adversary can reach and the spatial envelope within
which the mission must remain (\S\ref{sec:surfaces}); the capabilities assumed
of the adversary and of the detector (\S\ref{sec:adversary}); and the precise
meaning of mission completion (\S\ref{sec:mcr}). \S\ref{sec:formulation} then
defines the object that existing datasets do not provide and that the rest of
the paper is designed to supply.

\subsection{Attack Surfaces and Operating Volume}
\label{sec:surfaces}
We partition the UAV attack surface into two \emph{layers}. The \layernet{}
layer covers everything between nodes: the swarm mesh (multi-hop routing among
aircraft), the C2 link, and the intra-UAV bus. The \layerauto{} layer covers
everything inside the navigation and control stack: the GNSS receiver, sensor
fusion, the flight controller, and the mission outcome. A \emph{sublayer}
refines each layer into the concrete surfaces that a dataset can observe.
Fig.~\ref{fig:surfaces} places the six datasets on the causal chain.

We also fix here the single spatial quantity on which the rest of the paper
depends. A mission is flown inside a \emph{safe corridor}: a tube of lateral
half-width $m$, the \emph{corridor margin}, centred on the planned route, within
which the aircraft must remain for the flight to be considered nominal. The
corridor is not a construct of our analysis; it is the standard unit of
operational approval for uncrewed flight. Under the JARUS SORA methodology, the
operational volume is extended by a contingency volume whose lateral extent is
the sum of GNSS accuracy, position-holding error, map error, and the stopping
distance of the contingency manoeuvre~\cite{jarus_sora25}; for a small
multirotor these terms total roughly $7$\,m. We therefore report every certified
quantity over a family of margins $m\in\{2,5,10,20\}$\,m spanning that figure,
rather than at a single value. Leaving the corridor is the spatial failure mode;
\S\ref{sec:mcr} adds the temporal one.

\begin{figure}[t]
\centering
\resizebox{\columnwidth}{!}{%
\begin{tikzpicture}[font=\footnotesize,>={Latex[length=2mm]},
  node distance=2.5mm,
  n/.style={rounded corners=2pt,draw=netcol,fill=netcol!10,align=center,
    minimum height=6.5mm,inner sep=2.5pt,text width=17mm},
  a/.style={rounded corners=2pt,draw=autocol,fill=autocol!10,align=center,
    minimum height=6.5mm,inner sep=2.5pt,text width=17mm}]
\node[n] (mesh) {swarm mesh\\\scriptsize UAVIDS, UAV-CAS, DATAMUt};
\node[n,right=of mesh] (c2) {C2 link\\\scriptsize UAV Attack Dataset};
\node[n,right=of c2] (bus) {intra-bus\\\scriptsize HCRL UAVCAN};
\node[a,below=8mm of c2] (gnss) {GNSS + EKF\\\scriptsize state estimate};
\node[a,right=of gnss] (ctrl) {controller\\+ mission\\\scriptsize UAV-EW-Bench};
\node[a,left=of gnss] (nav) {navigation\\input $u(t)$};
% mesh feeds the navigation input directly (peer state corrections). Drawn to
% the node itself; an earlier version terminated on a coordinate expression
% that left the arrowhead pointing into empty space.
\draw[->,netcol] (mesh) -- (nav);
\draw[->,netcol] (c2) -- (gnss);
\draw[->,netcol] (bus) -- (ctrl.north);
\draw[->,autocol] (nav) -- (gnss);
\draw[->,autocol] (gnss) -- (ctrl);
\node[align=center,font=\scriptsize\itshape] at ($(mesh)!0.5!(bus)+(0,0.9)$)
  {network layer (attacks \emph{on the wire})};
\node[align=center,font=\scriptsize\itshape] at ($(nav)!0.5!(ctrl)-(0,0.95)$)
  {autonomy layer (attacks \emph{on the flight})};
\end{tikzpicture}%
}
\caption{Attack surfaces covered by the six datasets, arranged along the causal
chain. Network-layer events (top) propagate into the autonomy layer (bottom)
through the navigation input. The benchmark represents this downward link
explicitly rather than assuming it, and the theorem builds on that
representation.}
\label{fig:surfaces}
\end{figure}

\subsection{Adversary and Detector Model}
\label{sec:adversary}
We assume an adversary who can (i) inject, drop, delay, or forge traffic on one
or more network sublayers, and (ii) manipulate the RF environment of the GNSS
receiver. We further assume a network-layer \emph{detector} $D_\theta$ operating
at a tunable point $\theta$ that trades detection recall against false-positive
rate. For the multi-hop forwarding detector on which we build, $\theta$ is a
per-hop deviation tolerance in seconds, and a hop is flagged if its forwarding
delay exceeds $\theta$ or its next hop deviates from a scheduling oracle.
Lowering $\theta$ increases detections at the cost of more false positives;
raising it has the opposite effect. An \emph{evading} attacker keeps its
behaviour just below $\theta$ and is therefore undetected by construction. We
assume that the adversary knows $\theta$, which is the strongest reasonable
assumption and the one our worst-case analysis requires.

\subsection{Mission Completion: Spatial and Temporal Predicates}
\label{sec:mcr}
We adopt \emph{Mission Completion Rate} (MCR) as the operational metric, and
take $\MCR\ge0.90$ as a reference consistent with the airworthiness-security
posture of DO-326A/ED-202A~\cite{do326a}. Emerging guidance on AI assurance in
aviation~\cite{ed324} motivates treating a certified MCR floor as a first-class
safety property, and \S\ref{sec:theorem} provides such a floor.

A purely spatial definition of completion, under which the trajectory never
leaves the corridor of \S\ref{sec:problem}, is inadequate for the attack class
considered here. A time-delay adversary can keep the aircraft geometrically on
its route while causing it to arrive late, so that a delivery, inspection, or
search mission with a time-critical objective fails without any corridor
violation. We therefore define completion as the conjunction of a spatial and a
temporal predicate. Write $x(t)$ for the realised trajectory,
$x_{\mathrm{nom}}(t)$ for the nominal one, $m$ for the corridor margin in
metres, $t_{\mathrm{arr}}$ for the arrival time at the objective, $T$ for the
nominal mission duration in seconds, and $\kappa$ for the schedule tolerance
expressed as a fraction of $T$. Then
\begin{equation}
\label{eq:mcr}
\MCR \;=\; \Pr\Big[\underbrace{\sup_{t\in[0,T]}\Norm{x(t)-x_{\mathrm{nom}}(t)}\le m}_{\text{spatial}}
\;\wedge\;
\underbrace{t_{\mathrm{arr}}\le(1+\kappa)T}_{\text{temporal}}\Big].
\end{equation}
Where a quantity is computed under the spatial predicate alone, as for the
navigation anchor of \S\ref{sec:eval}, whose source records only geometric
completion, we label it \emph{Spatial MCR} and state this explicitly.

\subsubsection*{The probability space of Eq.~\eqref{eq:mcr}}
\label{sec:probspace}
Eq.~\eqref{eq:mcr} defines a probability, so its probability space must be
stated; leaving it implicit allows a population-level bound to be misread as a
per-flight one. Fix a detector operating point $\theta$. Then
$\delta(\theta)$ and $\Delta(\theta)$ are deterministic worst-case bounds rather
than random variables, so all the randomness in Eq.~\eqref{eq:mcr} comes
from the mission itself. Formally, let $(\Omega,\mathcal{F},\mathbb{P}_{\mathcal{M}})$
be a probability space in which a sample point $\omega\in\Omega$ is one mission:
a route, a nominal duration $T$, a corridor margin $m$, an initial state, a wind
and turbulence realisation, and the number $H$ of attacker-controlled hops on
the route it uses. $\mathcal{M}$ denotes the \emph{mission
distribution} and $\MCR=\MCR(\mathcal{M},\theta)$ is the
$\mathbb{P}_{\mathcal{M}}$-measure of the set of missions that satisfy both predicates.

Three consequences follow, which we state explicitly.

\begin{itemize}
\item \textbf{The guarantee is relative to $\mathcal{M}$.} Theorem~\ref{thm:main}
bounds $\MCR(\mathcal{M},\theta)$ for the $\mathcal{M}$ under which
$\delta(\theta)e^{LT_c}\le m$ is evaluated. A fleet whose missions are flown in
tighter corridors, or over routes with more attacker-controlled hops,
corresponds to a different $\mathcal{M}$ and inherits no bound from ours. In
particular
$\MCR(\mathcal{M}_{\text{benchmark}},\theta)\neq\MCR(\mathcal{M}_{\text{arbitrary}},\theta)$,
and the numbers we report concern the former.
\item \textbf{The dependence is confined to two parameters.} The mission enters
the bound only through the corridor margin $m$ and the hop count $H$. For this
reason we report every certified quantity over a family of margins rather than
at a single value; the family is a coarse marginalisation over the relevant part
of $\mathcal{M}$. An operator can therefore instantiate the bound for their own
$\mathcal{M}$ by reading off their own $m$ and $H$, without repeating the
analysis.
\item \textbf{A bound on a proportion is not a bound on a flight.} $\MCR\ge0.9$
permits one mission in ten to fail and says nothing about which missions fail. A
per-flight guarantee would require a bound conditional on $\omega$, which the
Gr\"onwall argument can provide only in the degenerate case where $m$ and $H$ are
fixed and known in advance. In that case the certificate is deterministic and
the probability redundant; otherwise the statement is necessarily
distributional.
\end{itemize}

The empirical mission distribution used throughout \S\ref{sec:eval} is that
induced by the navigation anchor's flight set, crossed with the margin family
$m\in\{2,5,10,20\}$\,m and the per-topology hop counts of
Table~\ref{tab:corpus}; we report it explicitly so that it need not be
reconstructed.

\subsubsection*{Sensitivity to the mission distribution}
\label{sec:missiondist}
The point above is often made but rarely measured, so we measure it here. The
navigation anchor flies $93{,}600$ missions over a full factorial design that
includes three mission profiles (cargo over mixed terrain, perimeter patrol,
search and rescue) and three receiver models; this is the axis along which the
mission distribution can be varied with everything else held fixed.
Table~\ref{tab:missiondist} reports MCR for each of the nine resulting
distributions.

The result depends on whether one conditions on the adversary, and this contrast
is the main finding. Marginalised over the full $0$--$40$\,dB jamming sweep, the
nine distributions are nearly indistinguishable: MCR lies in $[0.882,0.893]$, a
spread of $0.011$. Conditioned on the attack, they differ substantially: at
$J/S=35$\,dB the mission profile alone changes MCR by $0.214$, and at $30$\,dB
the receiver model alone changes it by $0.154$. Of most practical relevance, the
jamming power at which a distribution first falls below the $\MCR\ge0.90$
reference ranges from $18.1$ to $27.1$\,dB, a spread of $9$\,dB, or roughly a
factor of eight in adversary transmit power.

Two conclusions follow. First, averaging over adversary strength makes the
dependence appear negligible, yet an operator flies one profile against one
adversary and therefore operates at a point rather than at the average. Second,
on this evidence, a certified floor quoted without its mission distribution
carries an ambiguity of about $9$\,dB. We therefore state the distribution
wherever we quote a floor.

\begin{table}[t]
\centering
\caption{MCR under the composed defence for each of nine mission distributions,
marginalised over the $0$--$40$\,dB jamming sweep, with Wilson $95\%$ intervals.
The marginal spread is small; the spread conditional on attack strength
reaches $0.214$, and the $J/S$ at which MCR crosses $0.90$ spans $9$\,dB.
Provenance: simulation (physics-informed). Generated by
\code{experiments/mission\_distribution.py}.}
\label{tab:missiondist}
\scriptsize
\renewcommand{\arraystretch}{1.02}
\input{figures/block_paperD_tab_missiondist.tex}
\end{table}

Two conventions regarding $T$ prevent later ambiguity. First, $T$ is a duration
in seconds: the missions in our corpus last between $136$ and $246$\,s, and the
reference mission profile of \S\ref{sec:eval} is $70$\,s, so a schedule
tolerance of $\kappa=0.10$ corresponds to a permissible slip of several seconds.
Second, $T$ is not the horizon over which the trajectory-tube certificate
of \S\ref{sec:theorem} is propagated. That horizon, written $T_c$ and fixed at
$1$\,s throughout, is the interval between successive uncorrupted state updates,
after which the estimator re-anchors and the deviation bound restarts;
\S\ref{sec:theorem} makes this separation precise. Conflating the two would
either understate the mission or make the exponential bound vacuous.

The temporal predicate does not require a separate analysis. The same residual
budget $\Delta(\theta)$ that bounds the spatial perturbation also bounds the
schedule slip, because an evading attacker's total contribution to arrival delay
equals its accumulated undetected forwarding delay. The temporal predicate
therefore holds whenever
\begin{equation}
\label{eq:temporal_condition}
\Delta(\theta)\;\le\;\kappa\,T ,
\end{equation}
which by Lemma~\ref{lem:residual} is implied by $H\min(\theta,s)\le\kappa T$.
This adds a second edge to the certified window of \S\ref{sec:eval}, and it is
useful to note which of the two constraints is binding. At the operating point
we certify ($\theta=0.25$\,s, $H=2$) the induced slip is at most $0.5$\,s, so
for any mission longer than $5$\,s at a $10\%$ schedule tolerance the temporal
predicate is satisfied with ample margin and the spatial predicate binds.
The ordering reverses for a loosely tuned detector: at $\theta$ near the
contact-window slack $s=5$\,s the slip reaches $H s = 10$\,s, which exceeds a
$10\%$ tolerance on any mission shorter than $100$\,s. A short, time-critical
mission is thus governed by Eq.~\eqref{eq:temporal_condition}, and a long
mission in a tight corridor by the spatial bound; a purely spatial metric cannot
express this distinction.

\subsection{Formalising the Cross-Layer Pairing}
\label{sec:formulation}
Let a network attack occupy a time window $W_n=[t_0,t_1]$ on some network
sublayer, and let the induced navigation degradation occupy a window $W_a$ on
the autonomy layer. Existing datasets provide either $W_n$ (with packet labels)
or $W_a$ (with flight outcomes), but never the pair together with the mechanism
that links them. We formalise the missing object as a tuple
\begin{equation}
\label{eq:pairing}
P = \big(W_n,\; W_a,\; \theta,\; \Delta(\theta),\; \delta,\; \MCR\big),
\end{equation}
where $\theta$ is the detector operating point, $\Delta(\theta)$ the worst-case
residual effect of an attacker that evades $D_\theta$, $\delta$ the resulting
perturbation budget on the navigation input, and MCR the measured or certified
outcome.

The two halves of the paper operate on Eq.~\eqref{eq:pairing} in complementary
ways. The role of the benchmark is representational: it makes the complete
tuple $P$ recordable, storable, and auditable, whereas existing corpora record
only its first or second component. The role of the theorem is inferential: it
derives the final component, the mission outcome, from the third, the detector
operating point, through the intermediate quantities. This inference is the
chain
\begin{equation}
\label{eq:chain}
\theta \;\longmapsto\; \Delta(\theta) \;\longmapsto\; \delta(\theta)
\;\longmapsto\; \MCR \ge f\big(\delta(\theta)\big),
\end{equation}
read from left to right: the operating point $\theta$ bounds the residual
undetected attack $\Delta(\theta)$ (\S\ref{sec:interface}); the residual attack
maps to a bounded perturbation $\delta(\theta)$ on the navigation input
(\S\ref{sec:interface}); and a Lipschitz--Gr\"onwall trajectory-tube argument
turns that bounded input into a guaranteed MCR floor (\S\ref{sec:theorem}).
Individually, the three arrows rest on established results in detector
thresholding, worst-case attack analysis, and certified robustness; the novelty
lies in their composition.

The evidential value of $P$ depends on how $W_n$ and $W_a$ were linked, and we
therefore record this link explicitly in every record (\S\ref{sec:schema}).

\begin{definition}[Composed guarantee]
\label{def:composed}
Given a detector $D_\theta$ and navigation dynamics Eq.~\eqref{eq:ode}, a
\emph{composed guarantee} is a function $f$ such that for every attack $a$ with
$D_\theta(a)=\text{``clean''}$, the mission completes with rate
$\MCR(a)\ge f(\delta(\theta))$, where $\delta(\theta)$ bounds the
navigation-input perturbation any such $a$ can induce.
\end{definition}

% =====================================================================
\section{Related Work}
\label{sec:related}
% =====================================================================
Prior work divides along the two halves of this paper, into corpora on the data
side and bounds on the guarantee side. The same gap appears in both: each
community works within one layer and leaves the interface between layers to
assumption.

\paragraph{UAV Corpora and Dataset Construction}
Taxonomies of AI-enhanced UAV intrusion detection~\cite{islam2025uavids} report
a proliferation of datasets alongside a persistent weakness: general-purpose
network corpora lack UAV protocol semantics, while UAV-specific corpora each
cover a single sublayer. The consequence is measurable: models trained on
generic benchmarks misclassify a large fraction of GPS-spoofing events in real
autopilot logs~\cite{islam2025uavids}, which is direct evidence that the
network and navigation layers do not form a single undifferentiated feature
space. FANET surveys~\cite{tsao2022survey,ceviz2024survey} catalogue routing
attacks but stop at network-level effects. The individual corpora are
authoritative within their sublayers: \dataset{UAVIDS-2025}~\cite{uavids2025}
contributes $122{,}171$ NS-3 mesh flows, the \dataset{UAV Attack
Dataset}~\cite{whelan2020uav,whelan2020novelty} real PX4 flights under GPS
spoofing, jamming, and MAVLink ping-DoS, and \dataset{HCRL
UAVCAN}~\cite{hcrl_uavcan} bus-level flooding, fuzzing, and replay. None of them
extends to the flight.

Construction methodology affects validity as directly as scale does. Audits of
CIC-IDS2017 exposed labelling errors, flawed flow construction, and unrealistic
temporal separation, each of which inflated reported
performance~\cite{engelen2021troubleshooting,sharafaldin2018cicids}. A corpus
spanning two layers introduces a further failure mode, namely asserting a causal
link between a network event and a flight degradation that were never observed
together. The mandatory provenance flag of \S\ref{sec:schema} is designed to
bound this risk. Among released datasets the closest in ambition is
\dataset{UAV-CAS}~\cite{uavcas2026}, an AERPAW-calibrated swarm corpus and the
first to include collaborative multi-attack compositions; we ingest it as one of
our six sources. It nonetheless remains a network-layer benchmark, and its
authors list GNSS spoofing and jamming as future work. Where it asks how well a
collaborative mesh attack can be detected, we ask what that attack, whether
detected or missed, does to the aircraft.

\paragraph{Certified Robustness and Its Extension to Control}
Randomised smoothing~\cite{cohen2019smoothing,salman2019provably} yields a
probabilistic $\ell_2$ radius that scales to large models, whereas
Lipschitz-based certification~\cite{pauli2022training,zhang2022rethinking}
yields deterministic bounds that are known to be globally loose; neural
ODEs~\cite{chen2018neuralode} admit the Gr\"onwall-type
argument~\cite{yan2020robustnode} that we adopt. All of these bound the output
of a model. In the extension to sequential decisions, CROP certifies a
cumulative-reward lower bound under state perturbation~\cite{wu2022crop}, CARRL
certifies robust actions~\cite{everett2021carrl}, reachability tools bound
trajectories of neural-network-controlled systems~\cite{ivanov2019verisig,tran2020nnv},
and control-theoretic work establishes safety under sensor false-data
injection~\cite{li2022lqg}. CROP is the nearest analogue to a certified mission
floor, but it certifies a policy against a given perturbation ball and does not
address where that ball comes from.

\paragraph{Secure Estimation, Runtime Assurance, and Delay Attacks}
A parallel control-theoretic literature asks how much sensor corruption an
estimator can tolerate: exact recovery holds below a corruption
threshold~\cite{fawzi2014secure}, satisfiability-modulo-theory solvers make the
search practical~\cite{shoukry2017smt}, and recent surveys catalogue the
field~\cite{ding2021survey}. Such results bound estimation error given a limit
on how many sensors are compromised; they do not model a detector whose threshold determines
what an attacker can inject in the first place, which is the premise supplied by
Lemma~\ref{lem:residual}. Simplex and its
descendants~\cite{sha2001simplex,mehmood2022blackbox,cofer2020rta} guarantee
safety by switching from an unverified controller to a verified baseline when a
monitor is triggered. Our certified window defines the region in which the
high-performance policy may run, and thus supplies the quantitative boundary
that such architectures otherwise tune by hand. The destabilising effect of
injected delay is itself well studied~\cite{lou2019timedelay,etd2023uav,datamut2026},
but this literature establishes whether a delay attack is dangerous or
detectable and goes no further.

\paragraph{Composing Guarantees}
The closest prior art composes certificates within a single ML pipeline.
Sheikholeslami et al.~\cite{sheikholeslami2021provably} jointly certify a
classifier and a detection class, so that an input is provably flagged or
correctly classified; this is structurally our evade-or-bound dichotomy, but
confined to one layer. Yang et al.~\cite{yang2022sensing} decompose end-to-end
certified robustness across a two-stage sensing$\to$reasoning pipeline, the
strongest existing template for chain certification, but this too remains within
perception. Assume-guarantee and contract-based CPS design~\cite{pacti2024}
provides the compositional vocabulary but no ML-robustness certificates. No
prior result chains a network-intrusion-detection guarantee to a navigation
certificate. This paper closes that gap, which requires the corpus and the
theorem together.

% =====================================================================
\section{The Two-Layer Schema}
\label{sec:schema}
% =====================================================================
\subsection{Overview}
The schema addresses a single problem. The six datasets we combine describe
different phenomena, at different rates and without a common clock, so any
representation that forced them into one flat table would discard whatever did
not fit.

We therefore use three record types, arranged by increasing
interpretation (Fig.~\ref{fig:schema}).

\begin{enumerate}
\item An \textbf{Event} is a single observation as its source recorded it,
translated into common field names but never summarised. This layer is
deliberately minimal: it asserts nothing beyond what the source file contains.
\item An \textbf{AttackWindow} is a time interval during which one named
attack was active on one layer. This layer asserts that adversarial activity
occurred, and when.
\item A \textbf{CrossLayerPairing} joins one network window to one autonomy
window and asserts that the first caused the second. It is the only record in
the schema that makes a causal claim, and hence the only one that carries a
mandatory flag recording the strength of the evidence for that claim.
\end{enumerate}

The rest of this section describes each type in turn. The design
principle to carry forward is the last: causality is confined to one record
type, made explicit and graded, so that no analysis can rest silently on a
coupling that was never observed.

\subsection{The Three Record Types}
\textbf{Event} is one atomic observation (a mesh flow, a bus frame, a routing
hop, a telemetry sample or a complete flight) mapped to a small common core
(identity, layer/sublayer, relative time, endpoints and a unified attack label).
It also carries a \code{metrics} block holding the few numeric quantities shared
across sources (residual forwarding delay, position error, $J/S$,
mission-completed) and a \code{native} block that preserves every remaining
source column, so ingestion is lossless.

\textbf{AttackWindow} is a labelled attack interval within one scenario,
carrying the unified attack class and the \code{intensity} parameters
appropriate to that class. It is the unit of cross-layer alignment.

\textbf{CrossLayerPairing} realises the tuple $P$ of Eq.~\eqref{eq:pairing}. It
joins a network window to an autonomy window and records $\theta$, the residual
budget $\Delta(\theta)$, the perturbation $\delta$ with a named,
versioned mapping, the certificate hook and the empirical MCR. Its
\code{pairing\_basis} field records the strength of the evidence for the
record's causal claim and takes one of three values:
\code{measured\_same\_platform} (both windows recorded on one testbed; the
strongest evidence), \code{physics\_model} (coupled through a stated physical
model) or \code{synthetic\_alignment} (aligned by construction). A reviewer can
filter the benchmark by evidence strength, and no claim can silently
upgrade its own provenance.

\subsection{Rationale for the Schema}
Three design choices allow six heterogeneous sources to coexist without information
loss. Time is always relative to each source's capture start,
since the sources share no global clock, and cross-layer alignment is performed
at the window level; a mismatch between a $50$\,Hz telemetry stream and a
per-flow network record therefore never forces a lossy resampling. Attack
labels are unified into a common taxonomy, and the verbatim source label
is retained in \code{label\_native}. Every unmapped column is preserved,
which makes ingestion provably lossless.

Cross-layer alignment requires a stated rule and a stated error budget, since
it is the step at which a benchmark could silently create, or destroy, the
effect it reports. We anchor attack onsets under an affine map between the two
relative clocks and bound the residual misalignment by
$e_{\mathrm{align}}\le0.53$\,s for every released pairing. Because
$\Delta(\theta)$ is computed entirely within the network source, misalignment
can attach a pairing to the wrong autonomy window but cannot inflate or deflate
the quantity the certificate consumes. Appendix~\ref{app:alignment} gives the
map, the error terms and the rejection rule.

\begin{figure}[t]
\centering
\resizebox{\columnwidth}{!}{%
\begin{tikzpicture}[
  font=\footnotesize,
  box/.style={rounded corners=2pt,draw,align=center,minimum height=7mm,inner sep=3pt},
  net/.style={box,fill=netcol!12,draw=netcol},
  aut/.style={box,fill=autocol!12,draw=autocol},
  br/.style={box,fill=bridgecol!15,draw=bridgecol,thick},
  ev/.style={box,fill=black!4,minimum width=15mm},
  >={Latex[length=2mm]}]
\node[ev] (mesh) {mesh\\flow};
\node[ev,right=3mm of mesh] (c2) {C2\\frame};
\node[ev,right=3mm of c2] (bus) {bus\\frame};
\node[ev,right=3mm of bus] (hop) {route\\hop};
\node[ev,right=3mm of hop] (tel) {telem.\\sample};
\node[ev,right=3mm of tel] (fl) {flight};
\node[left=2mm of mesh] (evlab) {\textbf{Event}};
\node[net,below=7mm of c2,minimum width=30mm] (wn) {network AttackWindow $W_n$};
\node[aut,below=7mm of tel,minimum width=30mm] (wa) {autonomy AttackWindow $W_a$};
\node[br,below=8mm of $(wn)!0.5!(wa)$,minimum width=62mm] (pair)
  {\textbf{CrossLayerPairing} $P=(W_n,W_a,\theta,\Delta(\theta),\delta,\mathrm{MCR})$\\
   \scriptsize \code{pairing\_basis} $\in$ \{measured, physics-model, synthetic\}};
\draw[->] (mesh) -- (wn); \draw[->] (c2) -- (wn); \draw[->] (bus) -- (wn); \draw[->] (hop) -- (wn);
\draw[->] (tel) -- (wa); \draw[->] (fl) -- (wa);
\draw[->,bridgecol,thick] (wn) -- (pair);
\draw[->,bridgecol,thick] (wa) -- (pair);
\end{tikzpicture}%
}
\caption{Structure of the two-layer schema. Heterogeneous events from the six
datasets are folded into typed network and autonomy \emph{AttackWindows}; a
\emph{CrossLayerPairing} joins one of each and records the detector operating
point, the residual budget, the induced perturbation and the mission outcome,
tagged by how the coupling was established.}
\label{fig:schema}
\end{figure}

% =====================================================================
\section{The Corpus and Ingestion Adapters}
\label{sec:corpus}
% =====================================================================
Table~\ref{tab:corpus} summarises the six datasets and the layer each
anchors; Table~\ref{tab:ingest} reports what was actually ingested. Every
adapter emits schema-valid records and preserves unmapped source columns. We do
not re-host raw third-party data: the artifact ships integration scripts and
unified labels, and each source is fetched from its origin.

\begin{table*}[t]
\centering
\caption{The six-dataset corpus. Four network-layer sources and two
navigation-layer sources share one schema; only the last row reaches
the mission outcome. A single source observes both layers on one airframe. It
supplies same-platform calibration flights for $\gamma$ but, because it lacks
machine-readable attack intervals, no released \code{measured\_same\_platform}
pairings (\S\ref{sec:provenance_dist}).}
\label{tab:corpus}
\small
\renewcommand{\arraystretch}{1.2}
\begin{tabularx}{\textwidth}{@{}l l l l X@{}}
\toprule
\textbf{Dataset} & \textbf{Real/Sim} & \textbf{Layer $\cdot$ sublayer} & \textbf{Event kind} & \textbf{Unique role in the end-to-end benchmark}\\
\midrule
\dataset{UAVIDS-2025}~\cite{uavids2025} & Sim (NS-3.24) & \layernet{} $\cdot$ mesh & flow & Large-scale multi-hop mesh; blackhole/flooding/Sybil/wormhole at scale (122k flows).\\
\dataset{UAV-CAS}~\cite{uavcas2026} & Sim (digital twin) & \layernet{} $\cdot$ mesh & flow & AERPAW-calibrated swarm; collaborative multi-attack compositions (99k flows, 1024 configs).\\
\dataset{UAV Attack Dataset}~\cite{whelan2020uav} & Real testbed & \layernet{} $\cdot$ C2 (MAVLink) & telemetry & Real PX4 GPS spoof/jam + ping-DoS on the \emph{same} airframe family we model: the measured bridge.\\
\dataset{HCRL UAVCAN}~\cite{hcrl_uavcan} & Real testbed & \layernet{} $\cdot$ intra-bus & frame & On-board UAVCAN/DroneCAN flooding/fuzzy/replay, a third attack surface below the C2 link.\\
\dataset{DATAMUt} traces~\cite{datamut2026} & Sim (event) & \layernet{} $\cdot$ mesh (DTN) & hop & Time-window-graph forwarding detector with a \emph{sweepable} operating point $\theta$.\\
\dataset{UAV-EW-Bench-2026}~\ownbenchcite & Sim (physics-informed) & \layerauto{} $\cdot$ GNSS+ctrl & flight & The navigation anchor: full $0$--$40$\,dB $J/S$ sweep with per-flight MCR outcomes.\\
\bottomrule
\end{tabularx}
\end{table*}

\subsection{Adapter Design for Irregular Sources}
Each adapter performs lossless normalisation. Three merit description because
their sources are irregular. The \dataset{UAV Attack Dataset} adapter parses PX4
flight logs, computes horizontal position error against a benign reference (or,
where none exists, the flight's own pre-attack median) and emits the
same-platform calibration flights from which the rate $\gamma$ of
\S\ref{sec:interface} is estimated. Because the release contains no
machine-readable attack intervals, the adapter emits no
\code{measured\_same\_platform} pairing records (\S\ref{sec:provenance_dist}).
The \dataset{HCRL UAVCAN} adapter parses candump-style bus frames. The
per-scenario attack labels in the release are not machine-readable, so we derive
each scenario's class from the frames themselves: burst rate relative to the
normal bus rate, the fraction of distinct payloads and per-frame byte entropy
cleanly separate a flooding burst (one payload sent repeatedly above the bus
rate) from a fuzzy burst (near-random payloads) and from a replay burst (a small
set of captured payloads reused). Table~\ref{tab:hcrl} reports the result,
flagged for confirmation against the dataset's technical report. The
\dataset{UAV-CAS} adapter parses the flow statistics of the calibrated digital
twin, including its stringified list columns, and is verified against CSV files
produced by the original authors' generator.

\begin{table}[t]
\centering
\caption{The corpus as actually ingested. Counts come from schema-validated
adapter outputs; every record validates against the schema of
\S\ref{sec:schema}.}
\label{tab:ingest}
\small
\renewcommand{\arraystretch}{1.1}
\input{figures/block_paperC_tab_ingest.tex}
\end{table}

\begin{table}[t]
\centering
\caption{Attack classes derived from the frames of the \dataset{HCRL UAVCAN}
scenarios, inferred from burst rate, share of distinct payloads and byte
entropy; scenarios that exhibit more than one burst shape map to
\code{collaborative}.}
\label{tab:hcrl}
\small
\renewcommand{\arraystretch}{1.05}
\input{figures/block_paperC_tab_hcrl.tex}
\end{table}

% =====================================================================
\section{The Detector-Operating-Curve Harness}
\label{sec:harness}
% =====================================================================
A cross-layer guarantee requires the detector's behaviour as a function of its
operating point. We expose $\theta$ as a swept parameter and, for each value,
measure detection recall and false-positive rate over multiple seeds, three
network topologies, four routing policies and two attacker delay regimes.

The resulting operating curve (Fig.~\ref{fig:opcurve}) is the network-side input
to a composed guarantee: for any target false-positive budget, it gives the
largest undetected per-hop delay an evading attacker may inject, that is,
the residual budget $\Delta(\theta)$ of Eq.~\eqref{eq:pairing}. The knee in the
curve is structural. An evading attacker's per-hop delay is additive up to the
inter-node contact-window slack; beyond it, a missed contact window incurs a
full-period penalty.

The knee and the floor to which the curve descends reflect a structural
property of any store-and-forward schedule rather than a simulation artefact,
and both recur across all three topologies. Two observations carry over
directly to the guarantee. First, the floor height differs by topology because
it equals the probability that a malicious delay overshoots the slack and
exposes itself; this probability is higher in the medium delay regime
($U[1,10]$\,s) than in the low one ($U[1,7]$\,s), because a delay drawn from a
wider range is more likely to overshoot. Second, recall does not fall to zero
for any topology: beyond the knee the detector still catches the self-exposing
fraction, which is why an evading attacker's undetected budget saturates
instead of growing without bound. This is the empirical basis for
Lemma~\ref{lem:residual}.

\begin{figure}[t]
\centering
\begin{subfigure}{\columnwidth}
\centering
\begin{tikzpicture}
\begin{axis}[width=\columnwidth, height=3.9cm,
  title={\footnotesize low mode $U[1,7]$\,s},
  xlabel={$\theta$ (s)}, ylabel={recall},
  xmin=0, xmax=10, ymin=0, ymax=1.05,
  legend pos=north east, legend cell align=left, grid=both,
  grid style={black!12}, tick label style={font=\scriptsize},
  label style={font=\scriptsize}, legend style={font=\tiny}]
\input{figures/block_paperC_fig5_low.tex}
\end{axis}
\end{tikzpicture}
\end{subfigure}

\begin{subfigure}{\columnwidth}
\centering
\begin{tikzpicture}
\begin{axis}[width=\columnwidth, height=3.9cm,
  title={\footnotesize medium mode $U[1,10]$\,s},
  xlabel={$\theta$ (s)}, ylabel={recall},
  xmin=0, xmax=10, ymin=0, ymax=1.05,
  legend pos=north east, legend cell align=left, grid=both,
  grid style={black!12}, tick label style={font=\scriptsize},
  label style={font=\scriptsize}, legend style={font=\tiny}]
\input{figures/block_paperC_fig5_medium.tex}
\end{axis}
\end{tikzpicture}
\end{subfigure}
\caption{The detector operating curve across all three topologies (grid, ring,
double-ring), both delay regimes, and $8$ seeds $\times\,4$ routing policies.
Recall stays near its ceiling while $\theta$ remains below the $5$\,s
contact-window slack, then drops to a floor set by the missed-window
self-exposure penalty (a delay beyond the slack costs a full TWiG period and is
flagged at any $\theta$). The false-positive rate is zero throughout; the
measured benign residual ceiling is $0.178$\,s. The knee-and-floor shape is
universal; the floor height is the topology-specific self-exposure
probability.}
\label{fig:opcurve}
\end{figure}

\subsection{Detectors with a Learned Score}
We instantiate the harness on a threshold detector because its operating point
is a single, physically meaningful scalar, which makes the cross-layer sweep
unambiguous. The framework does not depend on this choice: any detector
exposing a tunable cutoff and a measurable score-to-magnitude profile reduces to
the same residual budget through an \emph{effective} operating point
$\theta_{\mathrm{eff}}$, defined in Appendix~\ref{app:learned}. Lemma~\ref{lem:residual} then
applies unchanged, with $\theta_{\mathrm{eff}}$ in place of $\theta$. We derive
this reduction but do not instantiate it; \S\ref{sec:limitations} records the
gap.

% =====================================================================
\section{The Interface Mapping $\theta \mapsto \delta(\theta)$}
\label{sec:interface}
% =====================================================================
The interface is where a network-security quantity becomes a navigation-security
quantity. It is the point at which the composition takes place and, as
\S\ref{sec:eval} shows, the quantity that determines whether a useful detector
setting is certifiable.

\subsection{Residual Budget of an Evading Attacker}
\begin{lemma}[Residual delay budget]
\label{lem:residual}
Let $D_\theta$ flag any forwarding hop whose injected delay exceeds $\theta$. An
attacker that evades $D_\theta$ injects at most $\theta$ of undetected delay per
malicious hop, but only up to the contact-window slack $s$: a per-hop delay
exceeding $s$ misses its forwarding window, incurs a full-period penalty
$P_{\!c}$ and is therefore flagged at any operating point. Undetected delay is
thus capped at $\min(\theta,s)$ per hop, and over a route with $H$
attacker-controlled hops the worst-case undetected cumulative delay is
\begin{equation}
\label{eq:delta_budget}
\Delta(\theta) \;\le\; H\,\min(\theta,\,s).
\end{equation}
\end{lemma}
\begin{proof}[Proof sketch]
Evasion requires a per-hop delay $\le\theta$. If $\theta\le s$, the delay fits
within the contact window and the additive term is $H\theta$. If $\theta>s$, any
delay in $(s,\theta]$ overshoots the window, waits a full period
$P_{\!c}\gg s$ and is flagged regardless of $\theta$; an evading attacker
therefore cannot use more than $s$ of invisible delay per hop, which gives the
$Hs$ cap. Summing over hops yields Eq.~\eqref{eq:delta_budget}.
\end{proof}

This cap is the knee: below the slack the residual budget grows linearly in
$\theta$, and above it the budget saturates at $Hs$, so the certified floor
plateaus instead of collapsing. We verify the cap directly on the
released harness. Over $15{,}392$ observed hops across three topologies, two
attack modes and eight seeds, no undetected malicious per-hop residual exceeds
$4.84$\,s ($s=5$\,s), and every larger delay is flagged as a missed window
($\ge55$\,s residual). That campaign attacker draws its delays independently of
$\theta$, so it tests the cap but not the bound below the slack. We therefore
also run a threshold-aware attacker, which adds the largest delay that keeps
each hop below $\theta$ after allowing for the benign forwarding jitter
($0.05$--$0.18$\,s). It is never flagged ($1{,}248$ runs over the same grid),
and its undetected delay never exceeds Eq.~\eqref{eq:delta_budget}. At
$\theta=0.25$\,s it realises a median of $73\%$ of the budget
$H\min(\theta,s)$ (at most $83\%$), the shortfall being the jitter allowance;
once the budget saturates it realises $99\%$, with no undetected hop above
$4.99$\,s. An attacker that injects $0.99\,\theta$ without the allowance is
flagged on every hop below the slack, and above it every delay misses its
window. This measurement replaces the classical additive budget
$H\theta+P_{\!c}\mathbf{1}[\theta>s]$ with the tighter, saturating bound of
Eq.~\eqref{eq:delta_budget}, so the benchmark refines the theorem in addition
to testing it.

\subsection{From Residual Delay to Input Perturbation}
A cumulative forwarding delay of $\Delta(\theta)$ seconds means that a
navigation correction arrives $\Delta(\theta)$ late, so the stack acts on state
stale by that amount. At bounded UAV speed $v_{\max}$, this staleness bounds the
state/position perturbation:
\begin{equation}
\label{eq:staleness}
\Norm{\delta u} \;\le\; v_{\max}\cdot \Delta(\theta).
\end{equation}
Equation~\eqref{eq:staleness} is the kinematic worst case: it assumes the
aircraft travels at top speed in the wrong direction throughout the staleness
interval. This is sound but pessimistic, because a real estimator does not fly
blind. When GNSS corrections are late or corrupted, the EKF down-weights them
and coasts on inertial data, so position error grows far more slowly. We
therefore tighten Eq.~\eqref{eq:staleness} to
\begin{equation}
\label{eq:gamma}
\Norm{\delta u}\;\le\;\gamma\cdot\Delta(\theta),\qquad \gamma\ll v_{\max},
\end{equation}
where $\gamma$ (metres of position error per second of degradation) is measured
on real GPS-spoofing and GPS-jamming flights from the same PX4 airframe family
the navigation model represents. Table~\ref{tab:whelan} reports the
calibration. A benign reference flight fixes the pre-attack noise floor, and the
jamming and spoofing flights give the post-onset growth rate under two explicit
choices of error signal: the raw receiver fix the attacker injects, and the
EKF-filtered position on which the controller acts. Because this choice is
itself a decision about certificate semantics, we expose it as an explicit
switch and report both; the certificate consumes the EKF value, which closes the
loop. The gap between the measured rate (at most $1.625$\,m/s,
\S\ref{sec:gammastability}) and $v_{\max}=15$\,m/s is the slack that determines
whether the detector's operating point is certifiable.

The calibration flights contain jamming and spoofing, not forwarding delay, so
transferring the rate from one to the other requires a physical argument. A
delayed peer correction leaves the estimator without fresh information for the
duration of the delay, and the EKF coasts on inertial data; physically, this is
the jamming regime, in which $\gamma$ is lowest (Table~\ref{tab:gamma}).
Spoofing is the more damaging regime, because counterfeit fixes actively pull
the estimate. We nevertheless evaluate the certificate at the spoofing
supremum. This is conservative for pure staleness, and it remains a
valid bound when a delayed report, fused against a state that has since moved,
passes the innovation gate and drags the estimate as a counterfeit fix would.

\subsection{Dependence of $\gamma$ on Operating Conditions}
\label{sec:gammastability}
Everything downstream consumes $\gamma$ as a single number, so whether it is in
fact a single number is the empirical question on which the paper most depends.
We therefore estimate $\gamma$ at every post-onset sample of the calibration
flights, $674$ in total, and record the conditions prevailing at each. The
answer is negative: $\gamma$ varies by $6.1\times$ between its median and its
supremum, and the variation is systematic rather than random.

\subsubsection{Choice of estimator}
The two available estimators of $\gamma$ disagree qualitatively, so
we state which one the certificate may use and why. Writing $e(t)$ for the
position error and $t_0$ for the attack onset,
\begin{equation}
\label{eq:gammaest}
\gamma_{\mathrm{cum}}(t)=\frac{e(t)}{t-t_0},
\qquad
\gamma_{\Delta}(t)=\frac{e(t)-e(t_0)}{t-t_0}.
\end{equation}
Onset is declared when the error crosses a detection threshold, so $e(t_0)$ is
non-zero: it is $0.64$\,m for jamming at the receiver and $9.27$\,m for
spoofing, where the injected fix has already been displaced a long way before
the threshold fires. Dividing this pre-existing offset by a small $t-t_0$ makes
$\gamma_{\mathrm{cum}}$ diverge, reaching $8.6$\,m/s below one second of
staleness. That value reflects the detector rather than the attack.
Lemma~\ref{lem:residual} bounds only the delay the attacker causes, so the
perturbation attributable to it is the error accrued since onset, and
$\gamma_{\Delta}$ is the estimator the composition may use. We report both. Had
we used $\gamma_{\mathrm{cum}}$, we would have concluded
$\gamma\approx8.6$\,m/s at short staleness, which exceeds
$\gamma_{\mathrm{req}}=6.14$\,m/s, and would have reported the framework as
unusable on the basis of a measurement artefact.

\subsubsection{Factors that influence $\gamma$}
Table~\ref{tab:gamma} gives the per-flight rates, and Table~\ref{tab:gammacamp}
the per-sample distribution binned by factor. Three dependences are visible,
each with a physical mechanism.

\emph{Attack type} changes $\gamma$ by $2.6\times$: spoofing increases error
faster than jamming. Jamming removes information, so the estimator gates out the
missing fix and dead-reckons at the inertial drift rate; spoofing injects
information that the estimator accepts, so error accrues at whatever rate the
adversary chooses. A plausible false fix therefore does more damage than an
absent one, which is why the delay class is certified at the spoofing rate
(\S\ref{sec:interface}).

The \emph{GNSS condition} shifts it in the expected direction: as the
number of satellites in view falls from $14$ to $4$, $\gamma$ increases.

\emph{Staleness} has the largest effect, and a non-monotonic one.
$\gamma_{\Delta}$ peaks at $1.625$\,m/s for $t-t_0\in[5,10)$\,s and decays to
$0.02$\,m/s beyond $40$\,s, because the error saturates while the denominator
continues to grow. A sound certificate must use the supremum rather than the
average.

A descriptive least-squares fit on standardised covariates explains
$R^{2}=0.79$ of the variation, with coefficients whose signs are physically
plausible for satellite count ($-0.22$), staleness ($-0.26$) and airspeed
($+0.12$). We report these as effect sizes and deliberately give no $p$-values:
samples within a flight are strongly autocorrelated and there are two attack
flights, so any nominal significance would be an artefact of the dependence
structure.

\begin{table}[t]
\centering
\caption{Per-flight rates $\gamma$: the coarse calibration obtained from a
single scalar per flight. The $2.6\times$ spread across attack types is visible
here; the finer structure requires Table~\ref{tab:gammacamp}. Generated by
\code{experiments/gamma\_stability.py}.}
\label{tab:gamma}
\small
\renewcommand{\arraystretch}{1.05}
\input{figures/block_paperD_tab_gamma.tex}
\end{table}

\begin{table}[t]
\centering
\caption{Per-sample distribution of the baseline-corrected rate
$\gamma_{\Delta}$, binned by factor, over $674$ post-onset samples. Counts are
shown because coverage is uneven: airspeed bins above $1$\,m/s are
sparse, since these are hover flights. Generated by
\code{experiments/gamma\_campaign.py}.}
\label{tab:gammacamp}
\scriptsize
\renewcommand{\arraystretch}{1.02}
\input{figures/block_paperD_tab_gamma_campaign.tex}
\end{table}

\subsubsection{Consequences for the certificate}
The certificate is evaluated at the supremum over the conditions it claims to
cover, which on this evidence is $\gamma=1.625$\,m/s; this is the only value of
$\gamma$ used in the remainder of the paper. It is $1.19\times$ the per-flight
secant rate of $1.365$\,m/s reported by a coarser calibration, so the secant is
mildly optimistic. With $H=2$, $\theta=0.25$\,s, $\hat L_{\mathrm{loc}}=1.181$
and $T_c=1$\,s, the tube radius is
$\rho=0.50\times1.625\times e^{1.181}\approx2.65$\,m, so the operating point is
certified at every corridor of $5$\,m or more but not at $2$\,m. At a
$10$\,m corridor the window edge is $\theta^{*}=0.94$\,s (the secant would give
$1.12$\,s).

\emph{Coverage.} The campaign varies attack type, GNSS condition and staleness
well, airspeed poorly (these are hover flights, with the $90$th percentile of
$\lVert v_{xy}\rVert$ below $0.65$\,m/s), and platform and flight mode not at
all: one airframe, hover only. A factor we could not vary is
reported as unvaried, never as having no effect. The $6.1\times$ spread is
accordingly a lower bound on the variability a deployment would encounter, and
condition~\ref{c:interface} of \S\ref{sec:scope} remains the weakest link in the
chain.

\emph{Temporal validity of the linear model.} Equation~\eqref{eq:gamma} is a
locally linear model, and its scope must be limited accordingly. While GNSS is
late rather than absent, the EKF continues to correct against its last accepted
fix and the position error grows approximately linearly with staleness. Once
the outage exceeds the filter's innovation-gating horizon, the estimator
propagates on inertial data alone, and uncorrected accelerometer bias and
gyroscope drift make the error grow quadratically with elapsed time rather than
linearly~\cite{kong2024px4ekf}. We therefore state the domain of the model
explicitly: Eq.~\eqref{eq:gamma} is used only for
$\Delta(\theta)\le\tau_{\mathrm{lin}}$, and beyond $\tau_{\mathrm{lin}}$ a
sound bound requires the second-order form
\begin{equation}
\label{eq:gamma_quad}
\Norm{\delta u}\;\le\;\gamma\,\Delta(\theta)\;+\;\tfrac12\,b_{\mathrm{IMU}}\,\big(\Delta(\theta)-\tau_{\mathrm{lin}}\big)_{+}^{2},
\end{equation}
with $b_{\mathrm{IMU}}$ the residual accelerometer-bias magnitude and
$(\cdot)_+$ the positive part. On the calibration flights the linear fit holds
to within its residual over the full observed degradation interval, giving
$\tau_{\mathrm{lin}}\approx2$\,s. The domain condition applies to the delay
$\Delta(\theta)=H\min(\theta,s)$, not to $\theta$ itself. At the operating point
$\theta=0.25$\,s, $\Delta=0.50$\,s lies well inside the linear domain. At the
upper edge of the window for a $10$\,m corridor ($\theta^{*}=0.94$\,s,
\S\ref{sec:eval}), $\Delta=1.89$\,s lies just inside it. For $\theta>1$\,s,
however, $\Delta(\theta)$ exceeds $\tau_{\mathrm{lin}}$, and this includes the
upper window edge for a $20$\,m corridor ($\theta^{*}=1.89$\,s,
$\Delta=3.78$\,s). There Eq.~\eqref{eq:gamma_quad} is the sound form, and the
floors and window edges reported in that range, which use the linear model, are
upper bounds rather than certified values; we do not have a measured
$b_{\mathrm{IMU}}$ with which to evaluate the quadratic term. No result at the
operating point depends on that range. The quadratic term grows further as
$\Delta(\theta)$ approaches $Hs=10$\,s, which is another reason why the certified
window is bounded above rather than open.

\begin{table}[t]
\centering
\caption{The measured delay-to-position rate $\gamma$ from three real PX4
flights (one benign reference, one GPS-jamming, one GPS-spoofing) in the
\dataset{UAV Attack Dataset}, self-referenced to the pre-attack median.
``noise'' is the $95$th percentile of pre-attack position error;
$\gamma_{\mathrm{lsq}}$ and $\gamma_{\mathrm{peak}}$ are the least-squares and
peak-secant growth rates after attack onset. The flights serve as same-platform
calibration data for $\gamma$; they are not released as pairing records
(\S\ref{sec:provenance_dist}).}
\label{tab:whelan}
\scriptsize
\renewcommand{\arraystretch}{1.1}
\input{figures/block_paperC_tab_whelan.tex}
\end{table}

% =====================================================================
\section{The Composition Theorem}
\label{sec:theorem}
% =====================================================================
We now assemble the components. \S\ref{sec:interface} bounded the
perturbation that an attacker who evades the detector can apply to the
navigation input. We now convert that input bound into a bound on the aircraft's
deviation from its planned route, and from there into a lower bound on the
fraction of missions that complete. The argument is short. It applies a
classical differential inequality to a quantity that had not previously been
measured.

\subsection{Navigation Dynamics and Certification Horizon}
We model the UAV navigation state $x(t)\in\Real^n$ as evolving under a (possibly
learned) drift $F$,
\begin{equation}
\label{eq:ode}
\dot{x}(t) = F\big(x(t), u(t)\big), \qquad t\in[0,T],
\end{equation}
where $u(t)$ is the sensor or measurement input and $T$ is the nominal mission
duration in seconds. A mission \emph{completes} in the conjunctive sense of
Eq.~\eqref{eq:mcr}: the realised trajectory remains within the safe corridor of
margin $m$ for all $t\in[0,T]$, and the objective is reached by the deadline
$(1+\kappa)T$. MCR is the probability of completion over the mission
distribution. An attack perturbs the input, $u\mapsto u+\delta u$; for the
time-delay class it also shifts the arrival time.

Deviation is not, however, propagated over all of $[0,T]$. A navigation filter
operates in closed loop: each time it accepts an uncorrupted measurement, it
re-anchors the state estimate, and any deviation accumulated since the previous
such update is discharged rather than carried forward. The interval between
consecutive uncorrupted updates is the \emph{certification horizon} $T_c$, over
which the trajectory tube must be propagated. We fix $T_c=1$\,s throughout,
matching the $1$\,Hz GNSS solution rate of the receivers in our corpus, so that
$T_c\ll T$ for every mission considered. The following assumption states this
requirement explicitly.

\begin{definition}[Re-anchoring interval]
\label{def:reanchor}
The navigation filter admits at least one measurement update not affected by the
attack in every window of length $T_c$. Each such update returns the attacked
estimate to the nominal one up to the filter's nominal estimation error, and the
corridor margin $m$ is defined net of that error. The deviation between the
attacked and nominal trajectories, measured from the update, therefore starts
at zero.
\end{definition}

The time-delay class satisfies Definition~\ref{def:reanchor} without further
assumptions: an evading forwarding attacker delays peer corrections by at most
$\Delta(\theta)$ per contact window but does not suppress them, so uncorrupted
updates continue to arrive. The same does not hold for a persistent-denial
adversary, a boundary recorded in \S\ref{sec:limitations}.

\begin{theorem}[Composed network-to-navigation guarantee]
\label{thm:main}
Let the navigation drift $F$ of Eq.~\eqref{eq:ode} be $L$-Lipschitz in its state
and its input over the operating region, let $T_c$ be a certification horizon satisfying
Definition~\ref{def:reanchor}, and let $m$ be the safe-corridor margin. Let
$D_\theta$ be a network detector with residual budget $\Delta(\theta)$
(Lemma~\ref{lem:residual}) and interface mapping $\delta(\theta)$
(Eq.~\eqref{eq:gamma} or the kinematic Eq.~\eqref{eq:staleness}). Then every attack
evading $D_\theta$ induces a trajectory-tube radius bounded by
\begin{equation}
\label{eq:tube}
\rho(\theta) \;=\; \delta(\theta)\,e^{L T_c},
\end{equation}
and the mission completes, in the conjunctive sense of Eq.~\eqref{eq:mcr},
whenever $\rho(\theta)\le m$ and $\Delta(\theta)\le\kappa T$; consequently
\begin{equation}
\label{eq:floor}
\MCR \;\ge\; f\big(\delta(\theta)\big) \;=\;
\Pr_{\text{missions}}\!\big[\,\delta(\theta)\,e^{LT_c}\le m
\;\wedge\; \Delta(\theta)\le\kappa T\,\big].
\end{equation}
\end{theorem}
\begin{proof}[Proof sketch]
Let $e(t)$ be the deviation between the attacked and nominal trajectories,
measured from the most recent re-anchoring instant. $L$-Lipschitzness of $F$
gives $\Norm{\dot e}\le L\Norm{e}+L\Norm{\delta u}$, and $e=0$ at each
re-anchoring by Definition~\ref{def:reanchor}. The comparison form of
Gr\"onwall's inequality~\cite{gronwall1919} then bounds $\Norm{e(T_c)}$ by
$\delta(\theta)(e^{LT_c}-1)\le\delta(\theta)e^{LT_c}$, which is
Eq.~\eqref{eq:tube}. By
Definition~\ref{def:reanchor} the same bound applies on every subsequent window,
so the supremum over $[0,T]$ also satisfies it and the spatial
predicate is discharged whenever $\rho(\theta)\le m$. The temporal predicate follows because
the only mechanism by which an evading time-delay attack displaces arrival is
its accumulated undetected delay, bounded by $\Delta(\theta)$ in
Lemma~\ref{lem:residual}. Taking probabilities over the mission distribution
gives Eq.~\eqref{eq:floor}. Appendix~\ref{app:proof} gives the full argument.
\end{proof}

\emph{On the Lipschitz hypothesis.} Theorem~\ref{thm:main} requires $F$ to be
$L$-Lipschitz over the operating region, and this qualifier has physical
content: the aircraft must remain in a smooth, locally linearisable regime, with
attitudes away from gimbal-lock singularities, airspeeds above stall, actuators
off their saturation limits, and no contact dynamics. These conditions hold
within the cruise and loiter envelopes of our corpus, and the constant
we use is measured inside that envelope rather than assumed globally
(\S\ref{sec:eval}). They can nevertheless fail. An aggressive evasive manoeuvre,
a stall-recovery transient, or actuator saturation under a strong gust can push
the dynamics into a regime in which no finite $L$ bounds $F$ on the relevant
set; there the certificate gives no bound, rather than a loose one. This
is also where the runtime-assurance interpretation of
\S\ref{sec:related} becomes operational: departure from the linearisable
envelope is the trigger that a Simplex-style monitor must use to hand control to
a verified fallback whose guarantee does not depend on $L$.

\begin{proposition}[Soundness]
\label{prop:sound}
The certified floor never overstates the achievable completion rate: for every
$\theta$, the true MCR of any $D_\theta$-evading attack satisfies
$\MCR\ge f(\delta(\theta))$.
\end{proposition}
\begin{proof}
Every step above is an inequality in the conservative direction: the Lipschitz
bound overestimates deviation growth, the input bound uses the worst-case
$\delta(\theta)$, and the corridor test is a sufficient, not a necessary, condition for
completion. Hence $f(\delta(\theta))$ lower-bounds the
true completion probability. The empirical validation of \S\ref{sec:eval}
confirms that the certified floor lies below the measured rate.
\end{proof}

Fig.~\ref{fig:chain} summarises the four stages of the argument and maps each
to the section that establishes it: the residual budget and the interface
mapping in \S\ref{sec:interface}, and the trajectory tube and corridor test in
this section.

Two consequences shape the remainder of the paper. First, Eq.~\eqref{eq:floor}
is monotone in $\theta$ below the slack $s$ and constant above it. Below $s$,
tightening the detector shrinks the residual budget and raises the certified
floor; above $s$, the budget saturates at $Hs$, so loosening the detector further
does not lower the floor. This quantifies both the navigation safety gained by
tightening network security and the point beyond which loosening it costs
nothing more. Second,
the exponential factor $e^{LT_c}$ in Eq.~\eqref{eq:tube} is the principal
limitation of the core certificate: for large $L$, or for a re-anchoring
interval $T_c$ stretched by a denial adversary, the deterministic bound becomes
loose. This motivates the complementary certificates of \S\ref{sec:certs}.

\begin{figure}[t]
\centering
\begin{tikzpicture}[
  font=\normalsize,
  stage/.style={rounded corners=3pt,draw,align=center,minimum height=8.5mm,
    text width=42mm,inner sep=4pt},
  net/.style={stage,fill=netcol!12,draw=netcol,thick},
  br/.style={stage,fill=bridgecol!15,draw=bridgecol,very thick},
  aut/.style={stage,fill=autocol!12,draw=autocol,thick},
  lbl/.style={font=\normalsize\bfseries,fill=white,inner sep=1.5pt},
  >={Latex[length=2.6mm]},
  node distance=8mm]
\node[net] (theta) {detector point $\theta$};
\node[net,below=of theta] (delta0) {residual budget $\Delta(\theta)$};
\node[br,below=of delta0] (delta) {perturbation $\delta(\theta)$};
\node[aut,below=of delta] (tube) {trajectory tube $\delta\,e^{LT_c}$};
\node[aut,below=of tube] (mcr) {certified floor $\MCR\ge f$};
\draw[->,thick] (theta) -- node[lbl,right] {Lemma~\ref{lem:residual}} (delta0);
\draw[->,thick] (delta0) -- node[lbl,right] {interface (\S\ref{sec:interface})} (delta);
\draw[->,thick] (delta) -- node[lbl,right] {Gr\"onwall} (tube);
\draw[->,thick] (tube) -- node[lbl,right] {corridor margin $m$} (mcr);
\begin{scope}[on background layer]
\node[fit=(theta)(delta0),draw=netcol!55,rounded corners,inner sep=6pt,
  label={[font=\small\bfseries,text=netcol]left:\rotatebox{90}{network side}}]{};
\node[fit=(tube)(mcr),draw=autocol!55,rounded corners,inner sep=6pt,
  label={[font=\small\bfseries,text=autocol]left:\rotatebox{90}{autonomy side}}]{};
\end{scope}
\end{tikzpicture}
\caption{The composition chain of Theorem~\ref{thm:main}. A detector operating
point passes from top to bottom through a residual-budget bound, an empirically
grounded interface mapping, a Gr\"onwall trajectory tube, and a corridor-margin
test, yielding a certified mission-completion floor. The upper two stages are
network-side quantities and the lower two autonomy-side; $\delta(\theta)$
links them.}
\label{fig:chain}
\end{figure}

% =====================================================================
\section{A Four-Certificate Architecture}
\label{sec:certs}
% =====================================================================
Theorem~\ref{thm:main} is built on the Lipschitz--Gr\"onwall certificate, but
that certificate alone does not suffice. Its bound is deterministic and
worst-case: it holds for any perturbation up to its radius, yet the radius
collapses as $LT_c$ grows, and it presumes a Lipschitz constant that can be
bounded tractably. Three regimes fall outside it, each covered by a
distinct certificate (Table~\ref{tab:certs}).

\begin{table}[t]
\centering
\caption{The four certificates cover disjoint regimes, and none subsumes
another. The Lipschitz--Gr\"onwall core provides the deterministic MCR floor;
the other three extend coverage where it is loose or silent.}
\label{tab:certs}
\small
\renewcommand{\arraystretch}{1.05}
\begin{tabularx}{\columnwidth}{@{}
  >{\hsize=0.82\hsize\raggedright\arraybackslash}X
  >{\hsize=0.70\hsize\raggedright\arraybackslash}X
  >{\hsize=1.48\hsize\raggedright\arraybackslash}X@{}}
\toprule
\textbf{Certificate} & \textbf{Guarantee type} & \textbf{Unique niche}\\
\midrule
Lipschitz--Gr\"onwall & deterministic, worst-case & The MCR floor: bounded input $\Rightarrow$ bounded tube $\Rightarrow$ bounded loss. Tight only for moderate $LT$.\\
Randomized smoothing~\cite{cohen2019smoothing} & probabilistic radius & Scales to large / non-Lipschitz-tractable models; certifies when the budget is known only statistically.\\
PAC-Bayes~\cite{viallard2021pacbayes} & distributional & Transfers the floor from the training mission distribution to unseen deployment regions.\\
MWU regret~\cite{arora2012mwu,balcan2015commitment} & online, adaptive & Bounds loss against an \emph{adaptive} jammer over time, a dimension no fixed-ball certificate addresses.\\
\bottomrule
\end{tabularx}
\end{table}

\paragraph{Independence of the Four Certificates}
The Lipschitz--Gr\"onwall radius defines a single deterministic ball.
Randomized smoothing certifies a probabilistic ball for models whose Lipschitz
constant is intractable; it trades determinism for scalability and so occupies
a distinct operating point rather than a special case. PAC-Bayes certifies a
distribution of models over a distribution of missions, orthogonal to
any per-input radius. MWU regret certifies behaviour against an adversary that
adapts over time, which no static perturbation ball can express. The
four form a cover rather than a hierarchy.

For concreteness, we state the two complementary bounds; the first is
instantiated numerically, and the second awaits one constant. Randomized smoothing~\cite{cohen2019smoothing} wraps the base model
in Gaussian noise of level $\sigma$ and certifies, with confidence $1-\alpha$
over $n$ samples, an $\ell_2$ radius
\begin{equation}
\label{eq:cohen}
R_{\mathrm{RS}} \;=\; \tfrac{\sigma}{2}\big(\Phi^{-1}(\underline{p_A})-\Phi^{-1}(\overline{p_B})\big),
\end{equation}
where $\Phi^{-1}$ is the inverse standard-normal CDF and
$\underline{p_A},\overline{p_B}$ are confidence bounds on the two largest class
probabilities. The complementary PAC-Bayes
bound~\cite{mcallester1999pac,viallard2021pacbayes} transfers a guarantee from
the training mission distribution to deployment. Given empirical risk
$R_{\mathrm{emp}}$ over $n_{\mathcal{M}}$ training missions and
prior-to-posterior divergence $\mathrm{KL}$, with probability $1-\beta$,
\begin{equation}
\label{eq:mcallester}
R_{\mathrm{true}} \;\le\; R_{\mathrm{emp}} + \sqrt{\tfrac{\mathrm{KL}+\ln(2\sqrt{n_{\mathcal{M}}}/\beta)}{2n_{\mathcal{M}}}}.
\end{equation}
Our engine evaluates Eq.~\eqref{eq:mcallester} using the trained model's measured
$R_{\mathrm{emp}}$; the only outstanding input is the numeric value of $\mathrm{KL}$,
which we mark as pending rather than substitute (Table~\ref{tab:constants}).


\emph{The four certificates on a single flight.} The four apply concurrently
rather than as alternatives, and each bounds a different object. On a
sortie through a corridor in which a compromised relay delays peer corrections
and a ground jammer varies its power, the Lipschitz--Gr\"onwall tube bounds the
trajectory deviation caused by the deterministic, bounded part of the
disturbance. It is the only one of the four to yield a per-instant geometric
envelope, so it alone discharges the corridor test of Eq.~\eqref{eq:mcr}.
Randomized smoothing bounds the stochastic part, and its claim concerns the
model's decision at a point. PAC-Bayes bounds the gap between the risk the model
exhibited on its training mission mix and its risk on this sortie; this is a
statement about the model, constant across the flight. MWU regret bounds
the defender's cumulative policy-switching loss against an adapting jammer, a
statement about the time axis that is vacuous at any single instant.

Because a trajectory, a decision, a model, and a policy sequence are distinct
objects, the four compose by conjunction rather than by selection: a deployment
takes its floor as the pointwise minimum
$\min\{f_{\mathrm{G}},f_{\mathrm{RS}},f_{\mathrm{PB}},f_{\mathrm{MWU}}\}$, and
the binding term indicates which assumption is closest to failing. In the
regime evaluated here (a bounded time-delay disturbance, in-distribution
missions and a non-adaptive adversary), the Gr\"onwall term binds and the other
three are slack; for this reason the main result of \S\ref{sec:eval} is a
Gr\"onwall curve. We nonetheless retain all four, since relaxing any one of
these conditions changes which term binds. The four also differ in maturity.
The Gr\"onwall and randomized-smoothing floors are instantiated numerically on
the reference model, and the MWU bound at $|S|=4$, whereas the PAC-Bayes bound
is reported in symbolic form, pending its numeric $\mathrm{KL}$
(Table~\ref{tab:constants}).

% =====================================================================
\section{Instantiation and Evaluation}
\label{sec:eval}
% =====================================================================
The composition theorem yields a usable guarantee only once its constants are
instantiated. This section supplies them and examines the result in three ways: whether the guarantee holds at a detector setting that would be
deployed in practice (\S\ref{sec:certfloor}), whether composing the layers
improves on either alone (\S\ref{sec:dominance}), and how much of the
construction rests on measurement rather than on modelling
(\S\ref{sec:provenance_dist}).

\subsection{The Navigation Anchor}
The autonomy layer supplies the result on which the construction is built.
Fig.~\ref{fig:mcr} shows MCR versus $J/S$ for four defence configurations,
computed from the $93{,}600$-flight per-flight table of the released benchmark.
The undefended baseline falls below the $0.90$ floor by $J/S\!\approx\!8$\,dB,
whereas the strongest configuration maintains it up to $J/S\!\approx\!25$\,dB.
Here the ranking of the defences is secondary: the autonomy layer produces a
graded, measurable outcome against which a network event can be paired.

\begin{figure}[t]
\centering
\begin{tikzpicture}
\begin{axis}[
  width=\columnwidth, height=4.4cm,
  xlabel={jamming-to-signal ratio $J/S$ (dB)},
  ylabel={Mission Completion Rate},
  xmin=0, xmax=40, ymin=0, ymax=1.03,
  legend pos=south west, legend cell align=left,
  grid=both, grid style={black!12},
  tick label style={font=\footnotesize},
  label style={font=\footnotesize},
  legend style={font=\scriptsize}]
\input{figures/block_paperC_fig4.tex}
\draw[black,dashed] (axis cs:0,0.90) -- (axis cs:40,0.90);
\node[anchor=west,font=\footnotesize] at (axis cs:0.5,0.850) {DO-326A floor 0.90};
\end{axis}
\end{tikzpicture}
\caption{Navigation anchor: MCR vs.\ $J/S$ on \dataset{UAV-EW-Bench-2026},
computed from the released $93{,}600$-flight table (Wilson $95\%$ intervals in
the artifact). Network attack windows are paired against this graded
outcome.}
\label{fig:mcr}
\end{figure}

\subsection{Reference Constants}
Every certified quantity in this section uses the measured local Lipschitz
constant rather than the global power-iteration estimate. This choice matters:
the local value is the larger of the two, so it widens the trajectory tube and
narrows the certified window; quoting the global value would make every bound
reported here more favourable. At $\theta{=}0.25$\,s the tube is $2.65$\,m, so
the operating point is certified for $m\ge5$\,m but not at $2$\,m; with the
global constant the tube would be $2.24$\,m, and the conclusion would not
change.

We instantiate the certificate on the reference navigation model, a
continuous-time graph model of the GNSS constellation and receiver. Global
power iteration gives a Lipschitz estimate $\hat L\approx1.01$; measuring the
constant along the hidden-state trajectories that the integrator visits on
operating-region inputs gives a larger local value
$\hat L_{\mathrm{loc}}\approx1.18$. We report the tighter local radius as the
defensible figure: with the certification horizon $T_c=1$\,s and an output
tolerance $\varepsilon_{\mathrm{out}}=0.5$,
$R=\varepsilon_{\mathrm{out}}e^{-\hat L_{\mathrm{loc}}T_c}\approx0.153$ (the
global estimate gives $0.182$). Since the local radius is smaller, the bound
is not loose in our favour, and the certified regime reflects the actual
sensitivity rather than an optimistic estimate. The
smoothing complement gives a certified $\ell_2$ radius $\approx0.44$ at
$\sigma=0.25$ ($\alpha=10^{-3}$, $n=200$; measured mean $0.446$ on the trained
checkpoint), and the MWU defender mixes four policies over $N$ rounds with regret
$R(N)\le\sqrt{N\ln|S|}$, $|S|=4$. Constants come from the trained reference
implementation on a transparent synthetic corpus and are held fixed;
Table~\ref{tab:constants} lists each with its provenance tag.

\begin{table}[t]
\centering
\caption{Instantiated constants and their provenance. The only numeric constant
still pending is the PAC--Bayes prior/posterior KL, which is flagged rather than
substituted.}
\label{tab:constants}
\small
\renewcommand{\arraystretch}{1.05}
\input{figures/block_paperA_tab2.tex}
\end{table}

\subsection{A Worked End-to-End Example}
\label{sec:worked}
To make Eq.~\eqref{eq:chain} concrete, we trace a single operating point through
it. Consider the grid topology, whose routes carry $H=2$ attacker-controlled
hops, and the detector setting $\theta=0.25$\,s. \emph{(1) Threshold to budget.}
Since $\theta<s=5$\,s, Lemma~\ref{lem:residual} gives $\Delta(\theta)\le H\theta=0.50$\,s
of undetected staleness; the threshold-aware attacker of \S\ref{sec:interface}
realises $0.36$\,s of it (median; maximum $0.41$\,s). \emph{(2) Budget to perturbation.} With the measured
rate $\gamma=1.625$\,m/s, Eq.~\eqref{eq:gamma} gives
$\Norm{\delta u}\le0.81$\,m; the kinematic bound would instead give $7.5$\,m.
\emph{(3) Perturbation to tube.} With $\hat L_{\mathrm{loc}}=1.181$ and
$T_c=1$\,s, the tube radius is $\rho=0.8125\,e^{1.181}\approx2.65$\,m.
\emph{(4) Tube to floor.} For any corridor margin $m\ge2.65$\,m, the mission is
certified to complete. Of the margin family $\{2,5,10,20\}$\,m, only the
tightest $2$\,m corridor fails, so the certified floor is $0.75$. Under the
kinematic bound, step (3) instead yields $\rho\approx24.4$\,m, which exceeds
every margin; for this reason the empirical calibration, not the theorem,
determines the framing.
Fig.~\ref{fig:budget} shows the measured undetected budget that feeds step (2)
and confirms the saturation predicted by Lemma~\ref{lem:residual}. Step (5), the
temporal predicate of Eq.~\eqref{eq:mcr}, is satisfied with a wide margin at this
operating point: the induced schedule slip is $\Delta(\theta)\le0.50$\,s,
below a $10\%$ tolerance for any mission longer than five seconds.

\emph{Choice of corridor margins.} The family $\{2,5,10,20\}$\,m brackets the
lateral buffers already required for operational approval
(\S\ref{sec:problem}), so the certified window can be compared with a
regulatory quantity rather than with a modelling choice. Relative to the SORA
contingency budget~\cite{jarus_sora25,easa_uspace2021}, $2$\,m is tighter than
any compliant lateral buffer and serves as a deliberately adversarial stress
case; $5$\,m lies within the default static budget; $10$\,m represents the full
contingency extent of a small multirotor once reaction distance is included;
and $20$\,m covers a fixed-wing platform, whose larger turn radius and cruise
speed enlarge the manoeuvre term. Read against these values,
Fig.~\ref{fig:sensitivity} turns the abstract question of whether $\theta$ is
certifiable into the operational one that an approval submission must
answer: whether the detector setting keeps the worst-case evading tube within
the contingency volume the operator has already declared.

\begin{figure}[t]
\centering
\begin{tikzpicture}
\begin{axis}[
  width=\columnwidth, height=3.9cm,
  xlabel={detector threshold $\theta$ (s)},
  ylabel={mean malicious delay (s)},
  xmin=0, xmax=10, ymin=0, ymax=26,
  legend cell align=left, legend columns=2,
  legend style={at={(0.5,1.06)}, anchor=south, font=\scriptsize, draw=black!30},
  grid=both, grid style={black!12},
  tick label style={font=\footnotesize},
  label style={font=\footnotesize}]
\input{figures/block_paperA_fig_budget.tex}
\end{axis}
\end{tikzpicture}
\caption{Measured residual budget underlying Lemma~\ref{lem:residual}
(scenario~1, mean over $8$ seeds $\times\,4$ policies). The undetected delay
that an evader passes through the detector (solid) saturates well below the
total injected delay (dashed). Beyond the contact-window slack, additional
delay produces missed windows that are flagged at any $\theta$, so the curve the
certificate consumes flattens rather than growing.}
\label{fig:budget}
\end{figure}

\subsection{The Certified Floor and Its Sensitivity to the Interface Mapping}
\label{sec:certfloor}
\emph{How the floor is computed.} Eq.~\eqref{eq:floor} is evaluated on the
mission distribution of \S\ref{sec:probspace}, reduced to the two parameters
through which a mission enters the bound. For each $\theta$, the hop count is
fixed at the value for the topology ($H=2$ for the grid), so
$\rho(\theta)=\gamma H\min(\theta,s)\,e^{LT_c}$ is a single number, and the
floor is the equally weighted fraction of the margin family $\{2,5,10,20\}$\,m
for which $\rho(\theta)\le m$. The floor therefore takes the values $0$,
$0.25$, $0.5$, $0.75$ and $1$; at $\theta=0.25$\,s, $\rho=2.65$\,m fits within
every margin except $2$\,m, giving a floor of $0.75$. The temporal conjunct
$\Delta(\theta)\le\kappa T$ is not evaluated numerically because the anchor
does not record mission durations. With $\kappa=0.1$ it holds for every mission
longer than $5$\,s at $\theta=0.25$\,s, and longer than $100$\,s once the budget
saturates at $Hs=10$\,s.

Fig.~\ref{fig:headline} reports the certified floor as a function of the
detector threshold for two interface mappings, since the choice between them
changes the conclusion. Under the kinematic worst case,
$\theta=0.25$\,s falls just outside the certified window; under the empirically
measured mapping it falls inside at every corridor margin of $5$\,m or more.
Fig.~\ref{fig:sensitivity} quantifies the available headroom: the mapping would
have to steepen beyond $\gamma_{\mathrm{req}}=6.14$\,m/s (at a $10$\,m corridor),
which is $3.8$ times the measured supremum rate, before the operating point left
the window.
We consider this the principal finding of the paper: whether cross-layer
certification is useful depends on a measurement of how attacks propagate into
state error, not on the proof machinery.

\emph{Effect of cruise flight.} Our calibration flights are low-speed hover
and loiter (\S\ref{sec:limitations}), and an obvious objection is
that $\gamma$ might scale with airspeed. An error-budget argument suggests
otherwise; we present it as reasoning, not as measurement. The quantity
$\gamma$ is the rate at which the position estimate degrades while a correction
is withheld, not the speed of the aircraft. During that interval the filter
continues to dead-reckon on inertial measurements, so error accrues at the rate
at which inertial propagation diverges from the true state. Hence the
measured $\gamma\le1.625$\,m/s lies almost an order of magnitude below the
kinematic $v_{\max}=15$\,m/s, which assumes the estimate does not propagate
at all. Of the terms in this divergence, only the cross-track error due to a
heading error $\psi_{\mathrm{err}}$ scales with airspeed, growing at
$v\,\psi_{\mathrm{err}}$; the remaining terms (accelerometer bias,
and attitude error acting on specific force) are set by sensor quality and are
largest under the continuous lateral acceleration of a loiter turn, not in
straight-and-level cruise. Taking the hover value as the airspeed-independent
floor, exceeding $\gamma_{\mathrm{req}}$ at $v=15$\,m/s would require
$\psi_{\mathrm{err}}\gtrsim(6.14-1.625)/15\approx0.30$\,rad, a sustained heading
error of about $17^\circ$. A PX4-class EKF keeps heading uncertainty within a
few degrees while GNSS is healthy and raises a yaw-innovation fault well before
that point, so cruise could invalidate the window only in a regime where the
estimator has already flagged itself as untrustworthy. We therefore expect the
window to hold in cruise, although only the cruise-inclusive campaign of
\S\ref{sec:limitations} can confirm this.

\begin{figure}[t]
\centering
\begin{tikzpicture}
\begin{semilogxaxis}[
  width=\columnwidth, height=4.3cm,
  xlabel={detector operating point $\theta$ (s, log scale)},
  ylabel={certified floor},
  xmin=0.05, xmax=10, ymin=0, ymax=1.05,
  legend cell align=left, legend columns=2,
  legend style={at={(0.5,1.06)}, anchor=south, font=\scriptsize, draw=black!30,
    inner sep=2pt, nodes={inner xsep=2pt},
    /tikz/every even column/.append style={column sep=0.4em}},
  grid=both, grid style={black!12},
  tick label style={font=\footnotesize},
  label style={font=\footnotesize}]
\input{figures/block_paperA_fig3.tex}
\end{semilogxaxis}
\end{tikzpicture}
\caption{Headline result: certified floor vs.\ detector operating point
$\theta$ for the time-delay class (grid topology, Gr\"onwall tube at the measured
local $L{=}1.181$, $H{=}2$, corridor-margin family $\{2,5,10,20\}$\,m). Under
the kinematic worst case (red), $\theta{=}0.25$\,s lies just outside the
certified window; under the rate measured on real GPS-spoofing and jamming
flights ($\gamma{=}1.625$\,m/s, the per-sample supremum), it lies inside at
every margin of $5$\,m or more (shaded: the window $[0.178,0.944]$\,s at
$m{=}10$\,m). At the tightest $2$\,m corridor the tube exceeds the margin
($2.65$ vs.\ $2.00$\,m). The floor is the equally weighted fraction of the
margin family whose corridor contains the tube. Both single-layer baselines
certify a zero floor against an evading attacker (black).}
\label{fig:headline}
\end{figure}

\begin{figure}[t]
\centering
\begin{tikzpicture}
\begin{loglogaxis}[
  width=\columnwidth, height=3.9cm,
  xlabel={corridor margin $m$ (m)},
  ylabel={window edge $\theta^{*}(m)$ (s)},
  xmin=1.8, xmax=22, ymin=0.015, ymax=4,
  xtick={2,5,10,20}, xticklabels={2,5,10,20},
  legend cell align=left, legend columns=1,
  legend style={at={(0.5,1.06)}, anchor=south, font=\scriptsize, draw=black!30},
  grid=both, grid style={black!12},
  tick label style={font=\footnotesize},
  label style={font=\footnotesize}]
\input{figures/block_paperA_fig4_sensitivity.tex}
\end{loglogaxis}
\end{tikzpicture}
\caption{Sensitivity of the window edge: the loosest detector setting
$\theta^{*}(m)=m/(\gamma e^{LT_c}H)$ whose worst-case evading tube still fits a
corridor of margin $m$, for the measured and the kinematic mapping. A mapping
certifies the operating point used in this paper iff its curve lies above the
$\theta{=}0.25$\,s line; the feasible region is bounded below by the measured
benign residual ceiling ($0.178$\,s, the FPR$=$0 floor). At $m{=}10$\,m the
measured rate meets the requirement $\gamma_{\mathrm{req}}{=}6.14$\,m/s with a
$3.8\times$ margin; the kinematic worst case fails at every margin.}
\label{fig:sensitivity}
\end{figure}

\subsection{Dominance over Single-Layer Baselines}
\label{sec:dominance}
We state the dominance claim precisely, rather than as blanket coverage of the
range. The composed guarantee dominates two baselines: (i) \emph{autonomy-only},
a navigation certificate with no detector, which must assume the full
unmitigated attack including the missed-window penalties (measured cumulative
delay up to $424$\,s) and therefore certifies a zero floor against an evading
attacker; and (ii) \emph{network-only}, a detector without a navigation
guarantee, under which any evading attack leaves the aircraft undefended, giving
a floor of zero by definition. The composed guarantee is the only
configuration with a nonzero certified floor. Under a per-seed Wilcoxon
signed-rank test with Holm correction, the composed floor exceeds the
autonomy-only floor at all $13$ operating points on the $\theta$ grid
($192$ paired runs per point, pooled over topologies, delay regimes, seeds and
policies), with every adjusted $p<10^{-13}$. The composition prevails because it alone uses the
detector to reduce the perturbation the certificate must absorb; through
the tightened budget of Lemma~\ref{lem:residual}, the floor then saturates
instead of collapsing beyond the knee.

The certified operating window is nonempty but bounded. Under the
empirically calibrated interface it is $[0.178,\,0.94]$\,s at a $10$\,m
corridor, and the operating point $\theta=0.25$\,s lies well inside it. Under
the conservative kinematic mapping it is empty at $10$\,m and narrows to
$[0.178,\,0.205]$\,s even at $20$\,m, which excludes $\theta=0.25$\,s. We claim
the former and report the latter as the worst-case fallback; we do not claim
coverage of the entire relaxation range.

\subsubsection*{Scope of the comparison}
The two baselines are ablations of our own construction, obtained by removing
one of the two layers. The comparison against them shows that the improvement
comes from the composition and is not an artifact of either layer alone. It
does not establish superiority over alternative approaches to assuring an
autonomous system, and we make no such claim.

We deliberately do not benchmark against reachability-based
verification~\cite{ivanov2019verisig,tran2020nnv}, secure state
estimation~\cite{fawzi2014secure,shoukry2017smt}, runtime assurance
architectures~\cite{sha2001simplex,mehmood2022blackbox,cofer2020rta}, or
certified reinforcement learning~\cite{wu2022crop,everett2021carrl}. This is a
limitation rather than a convenience: those methods answer different questions,
none of them takes a detector operating point as input, and there is therefore
no common axis on which to compare them. Reachability analysis bounds the states
a system can enter from a given input set but takes that set as given, and that
set is the quantity our network layer supplies. Secure state estimation
bounds the error given a limit on how many sensors are compromised; it does not model a
threshold that determines what an adversary may inject. Runtime assurance
guarantees safety by switching controllers when a monitor fires; our certified
window is a candidate specification for that monitor rather than a competitor
to it.

A fair comparison would require either extending those methods to accept a
detector threshold, or reducing ours to a fixed perturbation ball and comparing
tube tightness. The second option is feasible and worthwhile; it would establish
whether our Gr\"onwall tube is competitive with a modern reachability tool at
equal input radius. We have not done so, and until then the
contribution should be read as connecting two previously unconnected layers,
not as outperforming the state of the art in either.

\subsection{Certified versus Empirical Floors}
The certificate is sound iff the certified floor never exceeds the realised MCR
(Proposition~\ref{prop:sound}). At the anchor $J/S=20$\,dB, the strongest
defence attains an empirical $\MCR=0.938$ over $93{,}600$ real benchmark
flights, well above the certified target of $0.80$. The gap reflects the
conservativeness of the worst-case Gr\"onwall bound, and it narrows where the
empirically tightened interface replaces the kinematic worst case.

We keep two floors distinct throughout, since conflating them would overstate
the guarantee. The reference certificate targets a \emph{certified}
$\MCR\ge0.80$ at the anchor operating point, whereas the empirical benchmark
reports the DO-326A $0.90$ \emph{operational} floor holding further into the
band. The first is a conservative proof-derived bound and the second a
measured rate; every reported number is labelled accordingly.

\begin{table}[t]
\centering
\caption{A representative cross-layer pairing, in which a detector-evading
time-delay attack still induces measurable navigation staleness. Values are
from the released artifact (harness simulation; $\gamma$ at its per-sample
supremum; benchmark MCR).}
\label{tab:pairing}
\small
\renewcommand{\arraystretch}{1.05}
\input{figures/block_paperC_tab2.tex}
\end{table}

\subsection{A Worked Pairing for Each Evidence Class}
Table~\ref{tab:pairing} reports the pairing at a representative operating point.
The \code{pairing\_basis} flag is best explained through one example of each
class. \emph{Measured, same platform:} a pairing of this class links a real
network or RF event to the position error it caused, with both recorded on a
single airframe, so that its perturbation is measured rather than modelled. The
\dataset{UAV Attack Dataset} contains flights of this kind, but its release
lacks machine-readable attack intervals; the class is therefore empty in the
released artifact (\S\ref{sec:provenance_dist}), and the flights are used only
to calibrate $\gamma$.
\emph{Physics model:} a forwarding-attack window carries a residual budget in
seconds, and coupling it to a navigation window through the stated relation
Eq.~\eqref{eq:gamma} yields a \code{physics\_model} pairing. It rests on a
named, versioned physical relation that a reviewer can inspect and challenge,
and is the class the certificate consumes. \emph{Synthetic
alignment:} when a network window from one simulated source and an autonomy
window from another are aligned by construction, the pairing is tagged
\code{synthetic\_alignment}; this class offers maximal coverage with the
weakest causal claim, and the flag records as much. A user requiring a
strong claim restricts the data to the first class; one requiring breadth may
accept the third, with its provenance explicit.

\subsection{Provenance Composition of the Corpus}
\label{sec:provenance_dist}
The flag is useful only if its distribution is published, and we therefore report
it. Table~\ref{tab:provenance} gives the acquisition split across
the corpus: $200{,}610$ of $431{,}773$ ingested events ($46.5\%$) were
captured on real hardware, and the remainder were generated by three different
kinds of simulator. The split is favourable, but it answers the weaker of the
two questions a reviewer should ask, and conflating the two would overstate the
benchmark.

The stronger question is how many pairings rest on measurement, and here the
answer is more restrictive. A \code{measured\_same\_platform} pairing requires a
single source to have observed, on one airframe, both the network event and the
flight it degraded. Only one source in the corpus, the
\dataset{UAV Attack Dataset}, meets this requirement. Its three flights supply
the delay-to-position rate $\gamma$ that the certificate consumes
(Table~\ref{tab:whelan}; a real-corpus measurement obtained by
self-referencing each flight to its own pre-attack median), but the release does
not include machine-readable attack on/off timestamps. Without them, the
adapter cannot emit the \code{AttackWindow} boundaries that a pairing record
requires, and we decline to infer them. The number of released
\code{measured\_same\_platform} pairings is therefore zero at present: every
released pairing is either \code{physics\_model} (the network window coupled to
a navigation window through the stated relation Eq.~\eqref{eq:gamma}, and the
class the theorem consumes) or \code{synthetic\_alignment} (aligned by
construction under Appendix~\ref{app:alignment}, with the error budget
Eq.~\eqref{eq:align_err}).

We report this explicitly because exposing it is the purpose of the flag, and
because it identifies the binding limitation of the benchmark: the scarcity of
captures that observe both layers simultaneously, rather than the volume of
data, which is large. The headline guarantee does not depend on closing this
gap, since $\gamma$ is measured on real flights whether or not the window
boundaries are machine-readable; a reader filtering the artifact to
\code{measured\_same\_platform} will correctly retrieve nothing. Same-platform captures with published attack intervals, covering
flight regimes beyond hover, are consequently the most valuable extension of
the corpus.

\begin{table}[t]
\centering
\caption{Acquisition provenance across the ingested corpus. Counts come from
the schema-validated adapter outputs, and the split is machine-generated by
\code{experiments/provenance\_ledger.py} so the paper cannot drift from the
artifact. Acquisition is a property of the source; the strictly stronger
\code{pairing\_basis} distribution is discussed in the text.}
\label{tab:provenance}
\small
\renewcommand{\arraystretch}{1.05}
\input{figures/block_paperD_tab_provenance.tex}
\end{table}

% =====================================================================
\section{Discussion}
\label{sec:discussion}
% =====================================================================
\emph{What the composition offers an operator.} It converts a network-security
parameter into a navigation-safety figure. An operator who sets a detector's
false-positive budget currently reasons in packet statistics and has no means of
relating that choice to whether the fleet completes its missions.
Theorem~\ref{thm:main} supplies this link: each detector setting maps to a
certified mission floor, so the operator can instead ask how tight the detector
must be for a given corridor and read the answer from Fig.~\ref{fig:sensitivity}.

\emph{Why the interface, rather than the theorem, is decisive.} One might expect
the certificate to be the difficult part. Our results indicate otherwise: the
tube follows in a few steps from a classical inequality, and the Lipschitz
constant is measured rather than assumed. Whether a useful detector setting can
be certified depends on $\gamma$, and the order-of-magnitude gap between the
kinematic bound and the measured rate separates a vacuous certified window from
one that contains the operating point. Tightening cross-layer guarantees is
therefore primarily a measurement problem, which is why the benchmark and the
theorem appear in one paper.

\emph{Generality.} Although we instantiate the chain on a time-delay forwarding
attack and a GNSS navigation model, it is agnostic to both: any detector with a
monotone operating point and any perturbation-to-state mapping fit into
Lemma~\ref{lem:residual} and Eq.~\eqref{eq:gamma}, leaving the theorem and the
four-certificate cover unchanged. Extending the interface catalogue to
jamming-intensity-to-error and to bus-level attacks is direct future
work.

% =====================================================================
\section{Limitations and Future Work}
\label{sec:limitations}
% =====================================================================
The limitations below do not affect the stated contributions; each points to a
concrete extension.

\emph{The interface is calibrated rather than characterised.} This is the
paper's principal limitation and is condition~\ref{c:interface} of
\S\ref{sec:scope}. The rate $\gamma$ is measured
on three flights of a single PX4 airframe family, all in low-speed hover or
loiter. \S\ref{sec:gammastability} shows that even there it varies by $6.1\times$
once estimated per sample rather than per flight. The campaign covers attack type, GNSS condition and staleness well
but airspeed poorly, and it cannot vary two factors at all, namely platform (one
airframe) and flight mode (hover only); closing them requires
flights rather than analysis. The reported spread is therefore a lower
bound on the variability a deployment would meet.

The consequence can be quantified. If a campaign over these factors found
$\gamma>\gamma_{\mathrm{req}}=6.14$\,m/s at a $10$\,m corridor, an increase of
almost four times over the largest value we observe, the operating point would
leave the certified window, and the defensible claim would revert to the narrow
``certify-small, detect-large'' regime that the conservative kinematic bound
already guarantees. \S\ref{sec:eval} gives an error-budget argument that such an
increase would require a sustained heading error of roughly $17^\circ$; an
argument, however, is not a measurement. Either outcome would be informative. If
$\gamma$ proves stable across regimes, the result would be materially strengthened; if
it does not, the correct model is $\gamma=\gamma(\text{state},\text{input},\text{flight
mode},\text{attack type})$ and the certificate must be evaluated at the
supremum over the operational design domain. We regard determining which holds as
the most valuable extension of this work.

\emph{The certificate assumes that the filter re-anchors.} Theorem~\ref{thm:main}
propagates deviation over $T_c$ and relies on Definition~\ref{def:reanchor},
which requires one uncorrupted state update per window. A delay adversary
satisfies this by construction, postponing corrections without suppressing
them. An adversary that denies updates outright does not, and in
that case the certified radius grows as $e^{LT_d}$, where $T_d$ is the achieved
denial duration. Certifying against sustained denial requires either a link-layer bound
on outage length or a fallback whose guarantee does not depend on $L$, which is
the runtime-assurance direction of \S\ref{sec:related}. More generally, the
Gr\"onwall bound is conservative by design, which is why the architecture is a
four-certificate cover; fusing the deterministic, smoothing and PAC-Bayes floors
into a single tighter envelope remains open.

\emph{Evidence-class imbalance.} Only one of the six datasets couples a
network/RF event to a measured flight on one airframe, so it is the only source
from which \code{measured\_same\_platform} pairings could be built. Until its
attack intervals are published, it contributes same-platform calibration flights
for $\gamma$ and no pairing records (\S\ref{sec:provenance_dist}). This reflects
the public data available, and it is why the provenance flag is mandatory
rather than optional.

\emph{Calibration constants that remain data-gated.} The unit bridge that maps a
raw GNSS position error into receiver-level feature space, which serves the
RF-manipulation classes rather than the time-delay class targeted by the
theorem, is calibrated on a synthetic signal corpus pending access to real
in-phase and quadrature recordings. The PAC-Bayes bound
Eq.~\eqref{eq:mcallester} is connected to the trained model's measured empirical
risk but awaits its numerical KL divergence. The headline result depends
on neither, and both are tracked explicitly so that no synthetic constant is
reported as real.

\emph{Analytical navigation anchor and detector scope.} Completion outcomes come
from a physics-informed analytical Monte Carlo backend calibrated against
full-fidelity anchors, so we describe them as a calibrated analytical
benchmark. We also instantiate the harness on a single forwarding detector whose
operating point is one clean threshold, which makes the cross-layer sweep
unambiguous. Extending it to multi-dimensional operating points requires a
higher-dimensional sweep but no change to the schema.

\emph{Threats to validity.} \emph{Construct:} the provenance flag bounds the
risk that a recorded pairing misrepresents the causal link it claims, and a
reviewer can restrict any analysis to the measured class. \emph{Internal:} two
of our own measurement errors were caught because the provenance discipline made
them visible, and every figure is regenerated from committed data.
\emph{External:} four of six sources are simulations or calibrated twins, so
absolute detector-performance numbers are corpus-specific; what transfers is the
structure, a property of store-and-forward scheduling and
estimator behaviour rather than of any one simulator.

% =====================================================================
\section{Conclusion}
\label{sec:conclusion}
% =====================================================================
Attacks on a UAV's network and degradation of its flight have been studied
separately, because no existing resource allowed them to be studied together.
We built a benchmark in which this coupling can be represented and proved, on
its own objects, the conditional guarantee that the coupling enables: on the
benchmark side, a two-layer schema whose unit is the cross-layer pairing, six
validated adapters and a detector-operating-curve harness; on the guarantee
side, a residual budget that measurement tightens into the
saturating bound $\Delta(\theta)\le H\min(\theta,s)$, an interface rate
calibrated on real spoofing and jamming flights, and a composition theorem
chaining a detector's operating point to a certified mission-completion floor.

The evaluation supports three findings. First, among the configurations
compared, the composed guarantee is the only one that certifies a non-zero
mission floor against a detector-evading attacker, and it does so at every
operating point. Second, the empirically calibrated interface, an order of
magnitude tighter than the kinematic worst case, places a deployed detector's
operating point inside the certified window, making the guarantee operationally
relevant rather than only formal. Third, the binding constraint on cross-layer
certification is the interface, that is, the measurement of how attacks
propagate into state error, rather than the proof machinery.

The third finding also limits the result. Because the interface determines the
outcome, the result is only as strong as the calibration behind it, and ours
rests on three flights of one airframe in hover, over which $\gamma$ already
varies by $6.1\times$ with attack type, satellite visibility and staleness. What
we offer is therefore a conditional guarantee under an empirically calibrated
cross-layer interface. The conditions are stated in \S\ref{sec:scope}, the
calibration is released, and the experiment that would confirm or refute its
stability is specified in \S\ref{sec:limitations}. Under those conditions, UAV
network security and UAV autonomy security are joined by a single certificate.

%=====================================================================
% =====================================================================
\section{Ethics Considerations}
\label{sec:ethics}
% =====================================================================
This work involves no human subjects and no personally identifiable data. All
six datasets are public research corpora released for security research. We
redistribute none of the raw third-party data and ship only integration scripts
and unified labels, so that each source is fetched from its origin under its own
licence. No live aircraft, production network or third-party system was
attacked: every attack reported here is either replayed from a published capture
or generated in simulation on infrastructure we control.

On dual use, our harness quantifies how an attacker can remain below a
detector's threshold. We judge disclosure net-positive. Remaining below a
published threshold is the defining behaviour of an evading attacker and is
already assumed in the detector's own analysis; our central result bounds what
such an attacker achieves and shows that it saturates; and the practical
consequence is a defensive configuration procedure, namely reading the required
detector tightness from Fig.~\ref{fig:sensitivity} for a given corridor margin.
We disclose no previously unknown vulnerability in a deployed system, so no
coordinated-disclosure process applied; the authors of the underlying detector
are aware of this analysis of its operating point.

% =====================================================================
\section{Open Science and Artifact Availability}
\label{sec:openscience}
% =====================================================================
Everything needed to reproduce this paper is released as an open artifact: the
schema, the six ingestion adapters, the operating-curve harness, the certificate
engine, the statistical analysis and the figure generator. Every figure and
table in the paper is generated automatically from committed result files rather
than transcribed, so readers can regenerate them from the released data; each
carries a provenance tag (real corpus, simulation-derived, or stated modelling
parameter) in an artifact ledger. The harness is deterministic per seed,
and headline comparisons are released per seed with their
nonparametric tests. This discipline exposed two of our own measurement errors
before publication: a sweep harness that accumulated across operating points,
and a set of bus-attack labels that are not machine-readable per scenario, which
we therefore derive and flag rather than guess. We intend to submit the artifact for evaluation.
\ifanon
The repository is anonymised for review and will be de-anonymised on acceptance.
\else
The artifact is available at \url{https://robustuavs.ai}.
\fi

\ifanon\else

% =====================================================================
\section*{Acknowledgements}
We thank Keiwan Soltani for releasing the \dataset{DATAMUt} detector source and
for clarifying its operating point and contact-window semantics, which together
underpin the interface mapping of \S\ref{sec:interface}, and the maintainers of
the third-party datasets unified here.
\fi

\bibliographystyle{IEEEtran}
\bibliography{refs}

\appendix

\section{Background and Terminology}
\label{sec:background}
% =====================================================================
The composition proved in this paper spans two research communities with few
shared conventions, one concerned with packets and detector thresholds and the
other with trajectories and certified radii. This section fixes the single
system model to which both halves of the argument refer, namely how an
autonomous UAV closes its navigation loop, which of its interfaces an adversary
can reach, and what a robustness certificate does and does not promise, so that
the interface mapping of \S\ref{sec:interface} joins two precisely defined
quantities rather than two informal ones. Acronyms are defined at first use and
not tabulated; Table~\ref{tab:notation} is the single reference for
mathematical symbols.

\subsection{The UAV Navigation Loop}
An unmanned aerial vehicle (UAV, colloquially a drone) closes a control loop
around a \emph{state estimate}. An onboard computer fuses fixes from a Global
Navigation Satellite System (GNSS, of which the American GPS is one
constellation) receiver, inertial measurements from accelerometers and
gyroscopes (the IMU), and barometric height to obtain a best estimate of the
aircraft's position and velocity. In the PX4 autopilot family used throughout
this paper, fusion is performed by an extended Kalman filter (EKF), a standard
recursive estimator that weights each sensor according to its expected
noise~\cite{kong2024px4ekf}. The flight controller steers to keep the
estimate on the planned route.

Two properties of this loop underlie everything that follows. First, the loop is
only as good as its freshest trustworthy input: if GNSS is denied, the EKF
``coasts'' on inertial data and its position error grows with time. Second, the
loop is networked. Swarm members and the ground control station exchange
corrections and commands over three kinds of link: a command-and-control (C2) radio link,
typically using the Micro Air Vehicle Link (MAVLink) protocol; a multi-hop mesh
in which peers relay one another's traffic during periodic contact windows,
forming a delay-tolerant network (DTN); and an intra-vehicle bus, the UAV
Controller Area Network (UAVCAN, also DroneCAN), which connects the flight
controller to motor controllers and sensors. An attacker on the mesh can
therefore reach the navigation loop without touching the radio-frequency (RF)
front end.

\subsection{GNSS Spoofing and Jamming}
\emph{Electronic warfare (EW)} attacks the sensing physics. A \emph{jammer}
raises the jamming-to-signal power ratio ($J/S$, in decibels) until the
receiver loses satellite lock; a \emph{spoofer} broadcasts counterfeit satellite
signals so that the receiver locks onto a false
position~\cite{humphreys2008assessing,kerns2014capture}. Network attacks
instead target the messages. In a \emph{time-delay attack}, the class our
theorem covers, a compromised relay holds packets before forwarding them, so
that downstream consumers act on stale state~\cite{lou2019timedelay,datamut2026}.
Neither class needs to break cryptography; both degrade the control loop through
legitimate-looking inputs, which is why detection and robustness must be
composed rather than assumed.

\subsection{The Scope of a Robustness Certificate}
A \emph{certified} robustness statement differs from an empirical evaluation:
instead of reporting that test-set attacks were survived, it proves
that every perturbation within an explicitly stated budget is survived. The cost
is conservatism. Certified bounds are worst-case, so a certified floor typically
lies below what is observed empirically. We keep the two kinds of number
separate and label them throughout: certified floors come from proofs, empirical
rates from measured flights. We use two kinds of certificate: deterministic
Lipschitz-type bounds (if a model cannot amplify input changes by more than a
factor $L$, a bounded input change yields a bounded output change) and
probabilistic randomised-smoothing bounds (adding calibrated noise and taking a
majority vote yields a radius within which the decision provably cannot
change)~\cite{cohen2019smoothing}. The regulatory anchor for ``mission
survived'' is the airworthiness security process standard
DO-326A/ED-202A~\cite{do326a}, whose industry-acceptable completion floor of
$0.90$ appears in our figures as a horizontal reference line.


\begin{table}[t]
\centering
\caption{Mathematical notation. Each symbol is defined at first use; this
table serves as the single reference.}
\label{tab:notation}
\small
\renewcommand{\arraystretch}{1.05}
\begin{tabularx}{\columnwidth}{@{}lX@{}}
\toprule
$W_n$, $W_a$ & network / autonomy attack window\\
$P$ & cross-layer pairing tuple, Eq.~\eqref{eq:pairing}\\
$x(t)\in\Real^n$ & navigation state (position, velocity, \dots) at time $t$\\
$u(t)$ & sensor input to the navigation dynamics\\
$F$ & drift (right-hand side) of the navigation dynamics Eq.~\eqref{eq:ode}\\
$T$, $\kappa$ & nominal mission duration (s); schedule tolerance as a
fraction of $T$\\
$T_c$ & certification horizon: interval between uncorrupted state updates
over which the trajectory tube is propagated ($1$\,s)\\
$m$ & safe-corridor margin: allowed lateral deviation, metres\\
$D_\theta$, $\theta$ & network detector; its operating point
(tolerated per-hop residual delay, s)\\
$s$, $P_{\!c}$ & contact-window slack ($5$\,s); contact period ($60$\,s)\\
$H$ & number of attacker-controlled (malicious) hops on a route\\
$\Delta(\theta)$ & residual budget: worst undetected cumulative delay, s\\
$\delta u$ & perturbation of the navigation input caused by the attack\\
$\delta(\theta)$ & bound on $\Norm{\delta u}$ for an attack evading $D_\theta$, metres\\
$\gamma$ & delay-to-position rate: metres of error per second of degradation\\
$v_{\max}$ & kinematic worst-case speed bound ($15$\,m/s)\\
$L$, $\hat L$, $\hat L_{\mathrm{loc}}$ & (estimated / local) Lipschitz constant of $F$\\
$\rho(\theta)$ & trajectory-tube radius $\delta(\theta)e^{LT_c}$, metres\\
$\theta^{*}(m)$ & loosest $\theta$ whose tube fits margin $m$ (window edge)\\
$f(\cdot)$ & certified MCR floor as a function of the perturbation budget\\
$R$, $\sigma,\alpha,n$ & certified input radius; smoothing noise, confidence, samples\\
$R_{\mathrm{emp}}$, $\mathrm{KL}$ & empirical risk and prior-posterior divergence, Eq.~\eqref{eq:mcallester}\\
$n_{\mathcal{M}}$, $\beta$ & number of training missions; PAC-Bayes confidence parameter\\
$|S|$ & size of the defender's policy set in the MWU certificate ($4$)\\
\bottomrule
\end{tabularx}
\end{table}



\section{Detectors with a Learned Score}
\label{app:learned}
\label{sec:roc}
We instantiate the harness on a threshold detector because its operating point
is a single, physically meaningful scalar, which makes the cross-layer sweep
unambiguous. Since most deployed UAV intrusion detection relies on learned
classifiers, we show that the framework does not depend on this choice.

A learned detector emits a score $\sigma(a)\in[0,1]$ and flags an event when
$\sigma(a)\ge P_{\mathrm{th}}$; sweeping $P_{\mathrm{th}}$ traces the familiar
receiver operating characteristic, each cutoff giving a
$(\mathrm{FPR},\mathrm{TPR})$ pair just as each $\theta$ does in
Fig.~\ref{fig:opcurve}. The only object the composition consumes is the residual
budget (the largest attack effect that survives the detector), so the framework
generalises as soon as $P_{\mathrm{th}}$ maps to such a budget, which it does
through the detector's \emph{score-to-effect} profile. Let $g(\cdot)$ be the
monotonically increasing map from an attack's physical magnitude $\mu$ (here,
injected per-hop delay in seconds) to the detector's expected score,
$\sigma=g(\mu)$, measured by running the detector over labelled attacks of known
magnitude, which every corpus in \S\ref{sec:corpus} supports. An attacker
evading the cutoff must keep $g(\mu)<P_{\mathrm{th}}$, hence
\begin{equation}
\label{eq:pth}
\mu \;<\; g^{-1}(P_{\mathrm{th}}) \;=:\; \theta_{\mathrm{eff}}(P_{\mathrm{th}}),
\end{equation}
and Lemma~\ref{lem:residual} applies unchanged with $\theta$ replaced by the
effective operating point $\theta_{\mathrm{eff}}$. The threshold detector is the
special case $g=\mathrm{id}$, for which $\theta_{\mathrm{eff}}=\theta$. Two
properties of $g$ carry over into the guarantee. Monotonicity makes $g^{-1}$
well defined and the certified window an interval rather than a union; where a
classifier is non-monotone in $\mu$, one takes the conservative
$\theta_{\mathrm{eff}}=\sup\{\mu: g(\mu)<P_{\mathrm{th}}\}$, which is sound but
looser. Because a learned score is stochastic, $g$ is a
distribution rather than a function, so $\theta_{\mathrm{eff}}$ must be read at
an upper confidence bound on the score. This is the regime in which the
randomised-smoothing certificate of \S\ref{sec:certs}, rather than the
deterministic Gr\"onwall tube, is the appropriate downstream bound.

Any detector that exposes a tunable cutoff and a measurable score-to-magnitude
profile therefore fits without a schema change: the harness sweeps
$P_{\mathrm{th}}$ instead of $\theta$, records
$(\mathrm{FPR},\mathrm{TPR},\theta_{\mathrm{eff}})$ per cutoff, and the
downstream chain is unchanged. The framework does not require a detector to be
deterministic, only that its operating point reduce to a bound on the physical
attack magnitude that survives it.

% =====================================================================
\section{Aligning Two Relative Clocks}
\label{app:alignment}
Because the sources share no global clock, a \code{synthetic\_alignment} pairing
must state exactly how a network window and an autonomy window from different
captures are placed on a common axis, and how much error that placement can
introduce. Left implicit, this is the step at which a cross-layer benchmark
could silently create, or remove, the effect it reports.

Let the network source have relative clock $\tau_n\in[0,T_n]$ and the autonomy
source $\tau_a\in[0,T_a]$, each measured from its own capture start. An
alignment is an affine map
\begin{equation}
\label{eq:align}
\phi(\tau_n) \;=\; a\,\tau_n + b, \qquad a=\frac{T_a}{T_n},
\end{equation}
whose scale $a$ normalises the two capture durations and whose offset $b$ is
fixed by \emph{anchoring attack onsets}: we choose $b$ so that
$\phi(t_0^{\,n}) = t_0^{\,a}$, where $t_0^{\,n}$ is the first malicious event
labelled in $W_n$ and $t_0^{\,a}$ is the onset of degradation in $W_a$. The
onset is not judged by eye: $t_0^{\,a}$ is the first sample at which the
autonomy metric exceeds its pre-attack median by more than $3\hat\sigma$, where
$\hat\sigma$ is the robust (median-absolute-deviation) scale of the pre-attack
segment. Both anchors are therefore determined by the source labels and the
source's own noise floor rather than by the analyst.

Three error terms bound the residual misalignment $e_{\mathrm{align}}$:
\begin{equation}
\label{eq:align_err}
e_{\mathrm{align}} \;\le\; \underbrace{\tfrac{1}{2}(\Delta t_n+\Delta t_a)}_{\text{quantisation}}
\;+\;\underbrace{|a-1|\,T_a}_{\text{rate mismatch}}
\;+\;\underbrace{e_{\mathrm{onset}}}_{\text{anchor detection}},
\end{equation}
where $\Delta t_n,\Delta t_a$ are the sampling intervals of the two sources (a
per-flow network record and a $50$\,Hz telemetry stream give
$\tfrac12(\Delta t_n+\Delta t_a)\le\tfrac12(1.0+0.02)=0.51$\,s), the rate term
vanishes when both captures are replayed at their native rate ($a=1$, as for
every pairing we release), and $e_{\mathrm{onset}}\le\Delta t_a$ is the
one-sample latency of the $3\hat\sigma$ crossing. For the released
\code{synthetic\_alignment} pairings this gives
$e_{\mathrm{align}}\le0.53$\,s.

This bound prevents the alignment from biasing the reported staleness. The
quantity the certificate consumes, $\Delta(\theta)$, is computed entirely within
the network source from per-hop residual delays; the alignment determines only
the autonomy window with which a given network window is paired, never the
numerical value of $\Delta(\theta)$. Misalignment can therefore attach a pairing
to the wrong autonomy interval, but it cannot inflate or deflate $\Delta(\theta)$
itself. Its effect is confined to the empirical MCR column, and it is bounded
there too: since $e_{\mathrm{align}}\le0.53$\,s is two orders of magnitude
below the $60$\,s contact period that separates distinct attack windows, no
released pairing can be displaced onto a neighbouring window. Pairings whose
alignment error would exceed one tenth of the shorter window's duration are
rejected by the adapter rather than emitted with a caveat.

\section{Proof of Theorem~\ref{thm:main}}
\label{app:proof}
\begin{proof}[Proof of Theorem~\ref{thm:main}]
Fix a mission and a window between two consecutive uncorrupted updates, and
measure time from the first of them. Let $x_{\mathrm{nom}}(t)$ be the nominal
trajectory (the solution of Eq.~\eqref{eq:ode} under the attack-free input
$u(t)$) and $x_{\mathrm{att}}(t)$ the trajectory under $u(t)+\delta u(t)$, and
write $e(t)=x_{\mathrm{att}}(t)-x_{\mathrm{nom}}(t)$. Because the margin $m$ is
defined net of the filter's nominal estimation error
(Definition~\ref{def:reanchor}), the update leaves $e(0)=0$. From
Eq.~\eqref{eq:ode},
$\dot e(t)=F(x_{\mathrm{att}},u+\delta u)-F(x_{\mathrm{nom}},u)$, and
Lipschitz continuity of $F$ in state and input, together with
$\Norm{\delta u(t)}\le\delta(\theta)$, gives
\[
\frac{d^{+}}{dt}\Norm{e(t)}\;\le\;\Norm{\dot e(t)}\;\le\;L\Norm{e(t)}+L\,\delta(\theta),
\]
where $d^{+}/dt$ is the upper Dini derivative, needed only where $e(t)=0$. The
comparison function $r(t)=\delta(\theta)(e^{Lt}-1)$ solves
$\dot r=Lr+L\delta(\theta)$ with $r(0)=0=\Norm{e(0)}$, so the comparison
lemma~\cite{gronwall1919} gives $\Norm{e(t)}\le r(t)$ on the window. The same
result follows from the integral form of Gr\"onwall's inequality with forcing
term $g(t)=L\delta(\theta)t$:
$\Norm{e(t)}\le g(t)+\int_0^t g(\tau)Le^{L(t-\tau)}\,d\tau
=\int_0^t L\delta(\theta)e^{L(t-\tau)}\,d\tau$. At the certification horizon,
\begin{equation}
\label{eq:tube_derivation}
\Norm{e(T_c)}\;\le\;\delta(\theta)\big(e^{LT_c}-1\big)\;\le\;\delta(\theta)e^{LT_c},
\end{equation}
and we report the right-hand bound as Eq.~\eqref{eq:tube}. The slack between the
two bounds absorbs a second reading of re-anchoring. If an update leaves a
residual offset $\Norm{e(0)}\le\delta(\theta)$ instead of returning the state
exactly, and no further input perturbation acts until the next update, then
$\Norm{\dot e}\le L\Norm{e}$ and Gr\"onwall's inequality gives
$\Norm{e(T_c)}\le\delta(\theta)e^{LT_c}$ directly. Eq.~\eqref{eq:tube} is
therefore valid under either reading. By
Definition~\ref{def:reanchor} the deviation is reset at least once per window of
length $T_c$, so the same bound holds on each window and hence on the
supremum over $[0,T]$. Now suppose $\rho(\theta)\le m$. The nominal trajectory
lies in the safe corridor by construction, and every point of the attacked
trajectory is within $\rho(\theta)\le m$ of it by
Eq.~\eqref{eq:tube_derivation}; hence the attacked trajectory never exits the
corridor over $[0,T]$, discharging the spatial predicate of
Eq.~\eqref{eq:mcr}. For the temporal predicate, the only mechanism by
which a $D_\theta$-evading time-delay attack displaces arrival is the undetected
forwarding delay it accumulates, which Lemma~\ref{lem:residual} bounds by
$\Delta(\theta)$; hence $t_{\mathrm{arr}}\le T+\Delta(\theta)$, and
$\Delta(\theta)\le\kappa T$ gives $t_{\mathrm{arr}}\le(1+\kappa)T$ as required.
Since both predicates hold, the mission completes. Taking the probability of
$\{\rho(\theta)\le m\}\cap\{\Delta(\theta)\le\kappa T\}$ over the mission
distribution, the only source of randomness once $\theta$ is fixed (since
$\delta(\theta)$ and $\Delta(\theta)$ are deterministic bounds), yields
Eq.~\eqref{eq:floor}. Composing
$\theta\!\mapsto\!\Delta(\theta)$ (Lemma~\ref{lem:residual}),
$\Delta\!\mapsto\!\delta$ (\S\ref{sec:interface}), and $\delta\!\mapsto\!\MCR$
establishes the chain Eq.~\eqref{eq:chain}.
\end{proof}


\end{document}
